% content/c02-gpt/de.tex — Kapitel 2: Allgemeine probabilistische Theorien (deutsche Leitfassung)
%
% Quelle der Definitionen und zitierten Sätze: Plávala, arXiv:2103.07469v2
%   (Quellenschlüssel plavala-2103.07469v2; Dateien basic_v3.tex, what_v4.tex, CT_v3.tex,
%   QT_v3.tex, boxworld_v3.tex, channels_v3.tex, cones_v3.tex).
% Mathematik: evidence/derivations/c02/DERIVATION.md, geprüft in REVIEW.md daneben.
% Rahmung: docs/plan/BRIEF.md, Abschnitt „The physics approach“, Punkte 0 bis 3
%   (author-stipulated; nur zitiert und angeordnet, nicht erweitert).
%
% Schreibweise (Registry, TR-PLA-02 und TR-PLA-04): Zustände x, y; Effekte f, g, h.
% DERIVATION.md und IDEAS.md schreiben für dieselben Größen omega_i, e, pi_i; hier
% stehen x_i, f, lambda_i.  Die Koordinaten (b, a) eines Effekts sind die der Herleitung.

\chapter{Allgemeine probabilistische Theorien}\label{c02}

Dieses Kapitel legt den Rahmen fest, in dem der ganze Kurs arbeitet. Eine Theorie besteht
darin aus zwei Mengen: den Zuständen, in die sich ein System versetzen lässt, und den
Ja-Nein-Fragen, die man ihm danach stellen kann. Zu jedem Paar aus Zustand und Frage gehört
eine Wahrscheinlichkeit für die Antwort Ja. Drei Theorien rechnen wir vollständig durch: das
klassische Bit, das Quadrat-Bit und das Qubit. Am Ende steht eine Aufgabe, an der man die
drei auseinanderhalten kann: zwei Präparationen mit einer einzigen Messung unterscheiden.

Die Mathematik folgt der Einführung in allgemeine probabilistische Theorien von Plávala
\srcref{plavala-2103.07469v2}{„What are GPTs?“ und „State spaces and effect algebras“}.
Wie dort sind alle Zustandsräume dieses Kapitels endlichdimensional. Das ist eine Annahme
und keine Folgerung; sie steht als Bedingung (S5) in \defref{D3}.

\section{Was eine Theorie liefert}\label{c02.s1}

Der Kurs folgt einer Rahmung, die sein Autor festgelegt hat. Sie wird in diesem Kapitel
nicht begründet, sondern benutzt. Ihre ersten Sätze lauten:
% GAP-TJO: Die deutsche Fassung der drei folgenden Sätze ist eine Übersetzung aus
% docs/plan/BRIEF.md:44-46 und nicht von TJO geschrieben; der englische Wortlaut steht
% wörtlich in der englischen Ausgabe und in evidence/claims.json.
\begin{equation}\label{c02.eq.rahmen-liefert}\claim{c02.cl.rahmen-liefert}{author-stipulated}
  \begin{aligned}
    &\text{\tpcifenglish{A physical theory supplies data for}{Eine physikalische Theorie liefert Daten für}} \\
    &\text{\tpcifenglish{preparations + dynamics + measurement.}{Präparationen + Dynamik + Messung.}}
  \end{aligned}
\end{equation}
\begin{equation}\label{c02.eq.rahmen-probabilistisch}\claim{c02.cl.rahmen-probabilistisch}{author-stipulated}
  \begin{aligned}
    &\text{\tpcifenglish{All physical theories are probabilistic theories}{Alle physikalischen Theorien sind probabilistische Theorien}} \\
    &\text{\tpcifenglish{(I think of deterministic as a subset of probabilistic).}{(deterministische fasse ich als Teilmenge der probabilistischen auf).}}
  \end{aligned}
\end{equation}
\begin{equation}\label{c02.eq.rahmen-sprache}\claim{c02.cl.rahmen-sprache}{author-stipulated}
  \begin{aligned}
    &\text{\tpcifenglish{We use the language of preparations (= states)}{Wir benutzen die Sprache der Präparationen (= Zustände)}} \\
    &\text{\tpcifenglish{instead of “probability measures”.}{statt der „Wahrscheinlichkeitsmaße“.}}
  \end{aligned}
\end{equation}
Von den drei Zutaten in \eqref{c02.eq.rahmen-liefert} behandelt dieses Kapitel zwei: die
Präparationen und die Messungen. Die Dynamik kommt später im Kurs.

Was eine Präparation ist, legt \defref{D1} in Kapitel~\ref{c00} fest: eine Liste von
Anweisungen, nach der man ein System für ein Experiment herrichtet. Danach stellt man dem
System eine Ja-Nein-Frage. Die Theorie ordnet jedem Paar aus Präparation und Ja-Nein-Frage
eine Zahl zwischen $0$ und $1$ zu, die Wahrscheinlichkeit der Antwort Ja
\srcref{plavala-2103.07469v2}{„Effect algebra“}.
Abbildung~\ref{c02.fig.schema} zeigt einen Durchlauf eines solchen Experiments.

\begin{tpcfigure}{c02.fig.schema}
  \caption{Ein Durchlauf eines Experiments. Die Präparation versetzt das System in einen
    Zustand $x$; danach wird eine Ja-Nein-Frage gestellt. Die Theorie sagt die
    Wahrscheinlichkeit $f(x)$ der Antwort Ja voraus. Im ganzen Kapitel sind Zustände
    durchgezogen und als gefüllte Punkte gezeichnet, Effekte gestrichelt und als Quadrate.}
\end{tpcfigure}

Zwei Präparationen, die auf alle Ja-Nein-Fragen dieselben Wahrscheinlichkeiten liefern,
gelten nach \defref{D1} als derselbe Zustand. Ein Zustand ist also das, was an einer
Präparation für die Vorhersagen zählt.

\section{Zustandsräume}\label{c02.s2}

Die Menge aller Zustände eines Systems heißt sein Zustandsraum.

\tpcdef{D3}

In den Beispielen dieses Kapitels ist $\StSp$ eine Teilmenge von $\mathbb{R}^d$ mit
$d = 1$, $2$ oder $3$. Abgeschlossen und beschränkt bedeutet dann: $\StSp$ enthält seinen
Rand und reicht nicht ins Unendliche.

Warum konvex? Sind $x$ und $y$ Zustände und ist $\lambda$ eine Zahl zwischen $0$ und $1$,
so kann man würfeln: Mit Wahrscheinlichkeit $\lambda$ wird $x$ präpariert, mit
Wahrscheinlichkeit $1-\lambda$ wird $y$ präpariert. Auch das ist eine Präparation. Ihr
Zustand wird durch den Punkt $\lambda x + (1-\lambda) y$ beschrieben
\srcref{plavala-2103.07469v2}{„What are GPTs?“}.
Geometrisch liegt dieser Punkt auf der Strecke zwischen $x$ und $y$, und zwar umso näher bei
$x$, je größer $\lambda$ ist. Konvex heißt: Mit je zwei Punkten enthält $\StSp$ ihre ganze
Verbindungsstrecke.

\tpcdef{D4}

\begin{tpcfigure}{c02.fig.zustandsraeume}
  \demo{c02.demo.explorer}
  \caption{Die drei Zustandsräume dieses Kapitels: (a) eine Strecke, (b) ein Quadrat,
    (c) eine Kugel, hier mit ihrem Äquator gezeichnet. In jedem ist die Mischung
    $z = \tfrac14 x + \tfrac34 y$ zweier Zustände $x$ und $y$ eingetragen (offener Kreis);
    sie liegt auf der Verbindungsstrecke, dreimal so nah bei $y$ wie bei $x$. Die reinen
    Zustände sind in (a) die beiden Endpunkte $s_0$, $s_1$, in (b) die vier Ecken
    $s_{00}$, $s_{10}$, $s_{01}$, $s_{11}$ und in (c) alle Punkte der Kugeloberfläche.
    Das wird in Abschnitt~\ref{c02.s4} nachgerechnet.}
\end{tpcfigure}

Reine Zustände sind die Ecken des Zustandsraums im weitesten Sinn: Punkte, die auf keiner
Verbindungsstrecke zweier anderer Zustände im Inneren liegen. Alle übrigen Zustände lassen
sich aus ihnen mischen:

\begin{theorem}[Jeder Zustand ist eine Mischung reiner Zustände]\label{c02.thm.zerlegung}\claim[proof-omitted]{c02.cl.zerlegung}{source-backed}
  Sei $\StSp \subset V$ ein Zustandsraum in einem reellen Vektorraum $V$ der Dimension $d$.
  Dann lässt sich jeder Zustand $x \in \StSp$ als Konvexkombination von höchstens $d+1$
  reinen Zuständen schreiben: Es gibt reine Zustände $y_1, \dots, y_k$ mit $k \leq d+1$
  und Zahlen $\lambda_1, \dots, \lambda_k \geq 0$ mit $\lambda_1 + \dots + \lambda_k = 1$
  und
  \begin{equation}\label{c02.eq.zerlegung}
    x = \lambda_1 y_1 + \dots + \lambda_k y_k .
  \end{equation}
\end{theorem}

Das sind zwei Sätze der konvexen Geometrie. Plávala übernimmt sie ohne Beweis
\srcref{plavala-2103.07469v2}{„State space“}.
Im Quadrat kommt man demnach stets mit höchstens drei Ecken aus, in der Kugel mit höchstens
vier Punkten der Oberfläche. Eindeutig ist die Zerlegung im Allgemeinen nicht; das Quadrat
wird in Gleichung~\eqref{c02.eq.quadrat-mitte} das erste Gegenbeispiel liefern.

\section{Effekte}\label{c02.s3}

Eine Ja-Nein-Frage ordnet jedem Zustand $x$ eine Zahl $f(x)$ zu, die Wahrscheinlichkeit der
Antwort Ja. Für die Theorie ist sie also eine Funktion $f$ von $\StSp$ nach $[0,1]$. Nicht jede
solche Funktion kommt in Frage.

Um das zu sehen, betrachten wir die gewürfelte Präparation von eben auf zwei Arten. In
jedem Durchlauf entscheidet der Würfel, was präpariert wird: Mit Wahrscheinlichkeit
$\lambda$ zeigt er die Marke „$x$“, und es wird $x$ präpariert; mit Wahrscheinlichkeit
$1-\lambda$ zeigt er die Marke „$y$“, und es wird $y$ präpariert.

Erste Rechnung. Ein Durchlauf trägt mit Wahrscheinlichkeit $\lambda$ die Marke „$x$“ und
liefert dann mit Wahrscheinlichkeit $f(x)$ ein Ja. Mit Wahrscheinlichkeit $1-\lambda$ trägt
er die Marke „$y$“ und liefert dann mit Wahrscheinlichkeit $f(y)$ ein Ja. Zusammen ist die Wahrscheinlichkeit für Ja gleich
$\lambda f(x) + (1-\lambda) f(y)$.

Zweite Rechnung. Präpariert wurde der Zustand $\lambda x + (1-\lambda) y$. Die
Wahrscheinlichkeit für Ja ist $f(\lambda x + (1-\lambda) y)$.

Beide Rechnungen beschreiben dasselbe Experiment. Also muss für alle $x, y \in \StSp$ und
alle $\lambda \in [0,1]$ gelten:
\begin{equation}\label{c02.eq.affin}\claim{c02.cl.affin}{source-quoted}
  f(\lambda x + (1-\lambda) y) = \lambda f(x) + (1-\lambda) f(y) .
\end{equation}
Eine Funktion mit dieser Eigenschaft heißt \emph{affin}
\srcref{plavala-2103.07469v2}{„Effect algebra“}.

\begin{tpcfigure}{c02.fig.affin}
  \caption{Die Zustände mischen oder die Wahrscheinlichkeiten mischen. Waagerecht ist die Strecke zwischen zwei Zuständen aufgetragen, senkrecht die Wahrscheinlichkeit für Ja. (a) Eine affine Funktion: Ihr Graph ist eine Gerade. In der Mischung $z = \tfrac14 x + \tfrac34 y$ hat sie den Wert $\tfrac14\cdot\tfrac14 + \tfrac34\cdot\tfrac34 = \tfrac58$, dieselbe Mischung der beiden Randwerte. (b) Die Funktion $r \mapsto r^2$ auf $[-1,1]$: Das Mittel der Randwerte ist $1$, der Wert in der Mitte ist $0$. Eine solche Funktion kann keine Ja-Nein-Frage beschreiben.}
\end{tpcfigure}

\begin{greybox}[Zwei Annahmen hinter Gleichung~\eqref{c02.eq.affin}]
  Die Herleitung benutzt zwei Annahmen, die Plávala ausdrücklich nennt.
  \begin{enumerate}
    \item Mischungen lassen sich herstellen, indem man das Experiment oft wiederholt und
      einmal $x$, einmal $y$ präpariert.
    \item Es kommt dasselbe heraus, ob man die Zustände mischt oder die
      Wahrscheinlichkeiten. Anders gesagt: Die Wahrscheinlichkeit einer Antwort hängt nur
      vom Zustand ab und nicht davon, auf welchem Weg er hergestellt wurde.
  \end{enumerate}
  Aus der Konvexität von $\StSp$ allein folgt Gleichung~\eqref{c02.eq.affin} nicht. Auf
  der Strecke $[-1,1]$ etwa hat die Funktion $r \mapsto r^2$ Werte in $[0,1]$, aber in der
  Mitte zwischen $-1$ und $1$ den Wert $0$ und nicht den Mittelwert $1$ der beiden
  Randwerte.
\end{greybox}

\tpcdef{D2}

In Koordinaten sind affine Funktionen leicht zu erkennen\zoomref{c02.v1}.

\begin{lemma}[Affine Funktionen in Koordinaten]\label{c02.lem.koordinaten}\claim[proof-linked]{c02.cl.koordinaten}{derived}
  Sei $\StSp \subset \mathbb{R}^d$ ein Zustandsraum, der in keiner Geraden ($d = 2$), keiner
  Ebene ($d = 3$) und allgemein in keinem affinen Teilraum kleinerer Dimension enthalten
  ist; für $d = 1$ heißt das, dass $\StSp$ mehr als einen Punkt hat. Eine Funktion
  $f\colon \StSp \to \mathbb{R}$ erfüllt Gleichung~\eqref{c02.eq.affin} genau dann, wenn es
  eine Zahl $b$ und einen Vektor $\vb{a} \in \mathbb{R}^d$ gibt mit
  \begin{equation}\label{c02.eq.koordinaten}
    f(\vb{r}) = b + \vb{a}\cdot\vb{r}
    \qquad \text{\tpcifenglish{for all}{für alle}}\ \vb{r} \in \StSp .
  \end{equation}
  Die Zahl $b$ und der Vektor $\vb{a}$ sind durch $f$ eindeutig bestimmt.
\end{lemma}

\begin{vertiefung}[id=c02.v1,
    titel={Warum affine Funktionen diese Form haben},
    wozu={Der Kern benutzt Lemma~\ref{c02.lem.koordinaten} ohne Beweis. Hier steht er,
      zusammen mit der Regel für Mischungen aus mehr als zwei Zuständen.},
    voraussetzung={Kern bis Lemma~\ref{c02.lem.koordinaten}; lineare Gleichungssysteme.
      Der Rest des Kerns braucht diesen Block nicht.},
    zurueck={Auf einem Zustandsraum im $\mathbb{R}^d$ ist ein Effekt durch $d+1$ Zahlen
      festgelegt: $b$ und die Komponenten von $\vb{a}$.}]
  Gleichung~\eqref{c02.eq.affin} spricht von zwei Zuständen. Für $k$ Zustände
  $x_1, \dots, x_k$ und Gewichte $\lambda_i \geq 0$ mit $\lambda_1 + \dots + \lambda_k = 1$
  folgt daraus
  \begin{equation}\label{c02.eq.mehrfach}\claim{c02.cl.mehrfach}{derived}
    f(\lambda_1 x_1 + \dots + \lambda_k x_k) = \lambda_1 f(x_1) + \dots + \lambda_k f(x_k) .
  \end{equation}
  Das zeigt man schrittweise. Für $k = 2$ ist es Gleichung~\eqref{c02.eq.affin}. Ist
  $\lambda_k < 1$, so ist
  $\lambda_1 x_1 + \dots + \lambda_k x_k = (1-\lambda_k)\, y + \lambda_k x_k$ mit dem Zustand
  $y = \bigl(\lambda_1 x_1 + \dots + \lambda_{k-1} x_{k-1}\bigr)/(1-\lambda_k)$, der eine
  Konvexkombination von $k-1$ Zuständen ist. Gleichung~\eqref{c02.eq.affin} führt die
  Behauptung für $k$ auf die für $k-1$ zurück. Für $\lambda_k = 1$ ist nichts zu zeigen.

  \begin{proof}[Beweis von Lemma~\ref{c02.lem.koordinaten}]
    Jede Funktion der Form \eqref{c02.eq.koordinaten} ist affin, denn
    \begin{equation*}
      b + \vb{a}\cdot(\lambda\vb{r} + (1-\lambda)\vb{s})
      = \lambda(b + \vb{a}\cdot\vb{r}) + (1-\lambda)(b + \vb{a}\cdot\vb{s}) .
    \end{equation*}

    Sei umgekehrt $f$ affin. Nach Voraussetzung enthält $\StSp$ Punkte
    $\vb{v}_0, \dots, \vb{v}_d$, die in keinem affinen Teilraum kleinerer Dimension liegen.
    Dann hat das lineare Gleichungssystem $b + \vb{a}\cdot\vb{v}_i = f(\vb{v}_i)$,
    $i = 0, \dots, d$, für die $d+1$ Unbekannten $b$ und $\vb{a}$ genau eine Lösung. Wir
    setzen $g(\vb{r}) = b + \vb{a}\cdot\vb{r}$ und zeigen $f = g$ auf ganz $\StSp$.

    Jeder Punkt $\vb{r} \in \mathbb{R}^d$ lässt sich auf genau eine Weise als
    $\vb{r} = \alpha_0 \vb{v}_0 + \dots + \alpha_d \vb{v}_d$ mit
    $\alpha_0 + \dots + \alpha_d = 1$ schreiben; die $\alpha_i$ dürfen negativ sein. Wir
    bringen die Glieder mit negativem $\alpha_i$ auf die linke Seite. Sei $N$ die Menge
    dieser Indizes und $S = 1 + \sum_{i \in N} \abs{\alpha_i}$; die übrigen $\alpha_i$
    summieren sich dann ebenfalls zu $S$. Division durch $S$ ergibt
    \begin{equation*}
      \frac{1}{S}\,\vb{r} + \sum_{i \in N} \frac{\abs{\alpha_i}}{S}\,\vb{v}_i
      = \sum_{i \notin N} \frac{\alpha_i}{S}\,\vb{v}_i .
    \end{equation*}
    Für $\vb{r} \in \StSp$ stehen auf beiden Seiten Konvexkombinationen von Zuständen.
    Sowohl $f$ als auch $g$ erfüllen Gleichung~\eqref{c02.eq.mehrfach}, und sie stimmen in
    allen $\vb{v}_i$ überein. Wendet man beide auf diese Gleichung an und zieht die
    Ergebnisse voneinander ab, so bleibt $\bigl(f(\vb{r}) - g(\vb{r})\bigr)/S = 0$.
  \end{proof}
\end{vertiefung}

Nach Gleichung~\eqref{c02.eq.koordinaten} ist $f(\vb{r}) = b + \vb{a}\cdot\vb{r}$. Ist
$\vb{a} \neq 0$, also $f$ nicht konstant, so bilden die Punkte des $\mathbb{R}^d$, in denen
$b + \vb{a}\cdot\vb{r}$ denselben Wert hat, parallele Geraden ($d = 2$) oder parallele
Ebenen ($d = 3$), senkrecht zu $\vb{a}$; die Niveaumengen von $f$ sind ihre Schnitte mit
$\StSp$. Ist $\vb{a} = 0$, so ist $f$ konstant, und ganz $\StSp$ ist eine einzige
Niveaumenge. Ein Effekt ist eine affine Funktion, deren Werte auf ganz
$\StSp$ zwischen $0$ und $1$ bleiben (Abbildung~\ref{c02.fig.niveaulinien}). Affin ist nicht
dasselbe wie linear: Der konstante Anteil $b$ gehört dazu.

\begin{tpcfigure}{c02.fig.niveaulinien}
  \caption{Ein Effekt ist eine affine Funktion mit Werten in $[0,1]$. (a) Auf der Strecke
    ist sein Graph eine Gerade; sie ist durch die Werte $q_0$ und $q_1$ in den Endpunkten
    festgelegt, hier $\tfrac34$ und $\tfrac14$. Im Zustand $x$, der ein Viertel des Weges
    von $s_0$ nach $s_1$ liegt, liest man $f(x) = \tfrac58$ ab. (b) Auf dem Quadrat sind
    die Niveaulinien des Effekts $f(\vb{r}) = \tfrac12 + \tfrac14 r_x + \tfrac14 r_y$
    parallele Geraden senkrecht zu $\vb{a} = (\tfrac14, \tfrac14)$. Der Wert $0$ wird nur
    in der Ecke links unten angenommen, der Wert $1$ nur in der Ecke rechts oben; im
    Zustand $x = (\tfrac12, 0)$ ist $f(x) = \tfrac58$.}
\end{tpcfigure}

\tpcdef{D58}

Für jeden Effekt $f$ ist auch $\UnitEff{\StSp} - f$ ein Effekt, denn diese Funktion ist
affin und hat Werte in $[0,1]$:
\begin{equation}\label{c02.eq.komplement}\claim{c02.cl.komplement}{source-backed}
  \UnitEff{\StSp} - f \in \EffAlg{\StSp},
  \qquad f + (\UnitEff{\StSp} - f) = \UnitEff{\StSp} .
\end{equation}
Der Effekt $\UnitEff{\StSp} - f$ gehört zur Antwort Nein: Ist $f(x)$ die Wahrscheinlichkeit
für Ja, so ist $1 - f(x)$ die für Nein. Eine Ja-Nein-Frage ist damit dasselbe wie eine
Messung mit zwei Ergebnissen, beschrieben durch das Paar $f$, $\UnitEff{\StSp} - f$. Zwei
Effekte gibt es in jeder Theorie: $\UnitEff{\StSp}$, die Frage mit der sicheren Antwort Ja,
und die Nullfunktion $0$, die Frage mit der sicheren Antwort Nein.

Auch Effekte lassen sich mischen. Man würfelt, welche von zwei Fragen gestellt wird: mit
Wahrscheinlichkeit $\lambda$ die Frage $f$, sonst die Frage $g$. Die Wahrscheinlichkeit
für Ja ist dann $\lambda f(x) + (1-\lambda) g(x)$, und
\begin{equation}\label{c02.eq.effekt-mischung}\claim{c02.cl.effekt-mischung}{derived}
  \lambda f + (1-\lambda) g \in \EffAlg{\StSp}
  \qquad \text{\tpcifenglish{for all}{für alle}}\ f, g \in \EffAlg{\StSp},\ \lambda \in [0,1] ,
\end{equation}
denn diese Funktion ist wieder affin und hat Werte in $[0,1]$. Die Menge
$\EffAlg{\StSp}$ ist also selbst konvex, und sie hat Ecken:

\tpcdef{D74}

\section{Drei Theorien}\label{c02.s4}

\subsection{Das klassische Bit}\label{c02.s4.1}

\tpcdef{D11}

Das klassische Bit ist die klassische Theorie mit $n = 2$. Ein Zustand ist eine
Wahrscheinlichkeitsverteilung auf zwei Möglichkeiten, wie in Kapitel~\ref{c01}:
\begin{equation}\label{c02.eq.bit}\claim{c02.cl.bit-zustaende}{source-backed}
  \begin{aligned}
    \Simplex{2} &= \bigl\{ (\pw_0, \pw_1) \in \mathbb{R}^2 : \pw_0 \geq 0,\ \pw_1 \geq 0,\ \pw_0 + \pw_1 = 1 \bigr\}, \\
    s_0 &= (1,0), \qquad s_1 = (0,1) .
  \end{aligned}
\end{equation}
Das ist die Strecke zwischen $s_0$ und $s_1$ in Abbildung~\ref{c02.fig.zustandsraeume}(a);
jeder Zustand ist $x = \pw_0 s_0 + \pw_1 s_1$, und zwar auf genau eine Weise. Die
reinen Zustände sind die beiden Endpunkte $s_0$ und $s_1$
\srcref{plavala-2103.07469v2}{„Example: classical theory“}.

\begin{proposition}[Effekte des klassischen Bits]\label{c02.prop.bit}\claim[proof-included]{c02.cl.bit-effekte}{derived}
  Jeder Effekt $f$ auf $\Simplex{2}$ ist durch seine Werte $q_0 = f(s_0)$ und
  $q_1 = f(s_1)$ in den reinen Zuständen bestimmt:
  \begin{equation}\label{c02.eq.bit-effekt}
    f(\pw_0, \pw_1) = q_0\, \pw_0 + q_1\, \pw_1 .
  \end{equation}
  Umgekehrt ist das für jedes Paar $(q_0, q_1)$ mit $0 \leq q_0 \leq 1$ und
  $0 \leq q_1 \leq 1$ ein Effekt. Die Menge der Effekte ist also das Einheitsquadrat
  $\EffAlg{\Simplex{2}} = [0,1]^2$. Ihre Extremaleffekte sind die vier Ecken
  $(0,0)$, $(1,0)$, $(0,1)$ und $(1,1)$.
\end{proposition}

\begin{tpcfigure}{c02.fig.effektmengen}
  \caption{Die Effektmengen zweier Theorien. (a) Klassisches Bit: das Einheitsquadrat in
    den Koordinaten $q_0 = f(s_0)$, $q_1 = f(s_1)$. Der Punkt $f$ ist der Effekt aus
    Abbildung~\ref{c02.fig.niveaulinien}(a). (b) Quadrat-Bit, im Vorgriff auf
    Proposition~\ref{c02.prop.quadrat}: eine Doppelpyramide; nach oben ist $b$ aufgetragen, waagerecht $a_x$ und
    $a_y$. Die Extremaleffekte sind als Quadrate markiert, verdeckte Kanten gepunktet. Beim
    Bit gibt es vier Extremaleffekte zu zwei reinen Zuständen, beim Quadrat-Bit sechs zu
    vier reinen Zuständen.}
\end{tpcfigure}

\begin{proof}
  Gleichung~\eqref{c02.eq.bit-effekt} ist Gleichung~\eqref{c02.eq.affin} für
  $x = \pw_0 s_0 + \pw_1 s_1$. Umgekehrt ist die rechte Seite von
  \eqref{c02.eq.bit-effekt} affin, und ihr Wert ist ein gewichtetes Mittel von $q_0$ und
  $q_1$; er liegt genau dann für alle Zustände in $[0,1]$, wenn $q_0$ und $q_1$ dort
  liegen.

  Zu den Ecken: Liegt eine der beiden Zahlen, etwa $q_0$, echt zwischen $0$ und $1$, so
  kann man sie um einen kleinen Betrag $\varepsilon$ vergrößern und verkleinern, ohne das
  Quadrat zu verlassen; $(q_0, q_1)$ ist dann die Mitte zwischen
  $(q_0 + \varepsilon, q_1)$ und $(q_0 - \varepsilon, q_1)$ und kein Extremaleffekt. Sind
  dagegen $q_0$ und $q_1$ beide $0$ oder $1$ und ist $f = \lambda g + (1-\lambda) h$ mit
  Effekten $g$, $h$ und $0 < \lambda < 1$, so müssen $g$ und $h$ in jeder der beiden
  Koordinaten denselben Wert haben wie $f$: Ein Mittel zweier Zahlen aus $[0,1]$ ist nur
  dann $0$, wenn beide $0$ sind, und nur dann $1$, wenn beide $1$ sind.
\end{proof}

Die vier Extremaleffekte haben eine einfache Bedeutung. $(1,0)$ fragt: „Liegt $s_0$ vor?“
$(0,1)$ fragt: „Liegt $s_1$ vor?“ $(1,1)$ ist der Einseffekt und $(0,0)$ die Nullfunktion.
In den reinen Zuständen sind ihre Antworten sicher: Die Wahrscheinlichkeiten sind $0$ oder
$1$. Die übrigen Effekte sind Mischungen dieser vier, also Fragen, bei denen zusätzlich
gewürfelt wird. Abbildung~\ref{c02.fig.effektmengen}(a) zeigt die Effektmenge.

Zum Vergleich mit den beiden anderen Theorien legen wir den Nullpunkt in die Mitte der
Strecke. Mit der Koordinate $r = \pw_0 - \pw_1 \in [-1,1]$ ist
$\pw_0 = \tfrac12(1+r)$ und $\pw_1 = \tfrac12(1-r)$, und aus
Gleichung~\eqref{c02.eq.bit-effekt} wird
\begin{equation}\label{c02.eq.bit-zentriert}\claim{c02.cl.bit-zentriert}{derived}
  \begin{aligned}
    f(r) &= b + a\, r, \qquad b = \tfrac12(q_0 + q_1), \quad a = \tfrac12(q_0 - q_1), \\
    \abs{a} &\leq \min(b,\, 1-b) .
  \end{aligned}
\end{equation}
Die Ungleichung sagt dasselbe wie $0 \leq q_0 \leq 1$ und $0 \leq q_1 \leq 1$, denn
$q_0 = b + a$ und $q_1 = b - a$ sind der größte und der kleinste Wert von $f$, wenn
$a \geq 0$ ist, und umgekehrt, wenn $a < 0$ ist.

\subsection{Das Quadrat-Bit}\label{c02.s4.2}

Die zweite Theorie hat ein Quadrat als Zustandsraum. Plávala führt sie unter dem Namen
„boxworld“ ein: Ein Zustand beschreibt dort einen Kasten mit einem Eingabebit und einem
Ausgabebit, und die vier Ecken $s_{ij}$ sind die Kästen, die bei Eingabe $0$ die Ausgabe
$i$ und bei Eingabe $1$ die Ausgabe $j$ liefern
\srcref{plavala-2103.07469v2}{„Example: boxworld theory“}. In Koordinaten, deren
Nullpunkt in der Mitte liegt, ist
\begin{equation}\label{c02.eq.quadrat}\claim{c02.cl.quadrat-zustaende}{source-backed}
  \begin{aligned}
    \StSp_{\mathrm{Q}} &= [-1,1]^2 = \bigl\{ (r_x, r_y) : \abs{r_x} \leq 1,\ \abs{r_y} \leq 1 \bigr\}, \\
    s_{ij} &= (2i-1,\ 2j-1), \qquad i, j \in \{0,1\} .
  \end{aligned}
\end{equation}
Die reinen Zustände sind die vier Ecken (Abbildung~\ref{c02.fig.zustandsraeume}(b)). Anders
als beim Bit ist die Zerlegung in reine Zustände nicht eindeutig. Die Mitte des Quadrats
ist auf zwei Arten eine Mischung zweier Ecken:
\begin{equation}\label{c02.eq.quadrat-mitte}\claim{c02.cl.quadrat-mitte}{source-backed}
  \tfrac12\bigl(s_{00} + s_{11}\bigr) = \tfrac12\bigl(s_{10} + s_{01}\bigr) = (0,0) .
\end{equation}
In einem Simplex wäre das nicht möglich: Dort sind die Gewichte, mit denen ein Zustand aus
den reinen Zuständen gemischt ist, eindeutig bestimmt. Nach \defref{D11} ist das Quadrat-Bit
deshalb keine klassische Theorie.

\begin{tpcfigure}{c02.fig.quadrat}
  \caption{Das Quadrat-Bit. (a) Die Mitte des Quadrats liegt auf beiden Diagonalen; sie
    ist die Mitte zwischen $s_{00}$ und $s_{11}$ und ebenso die Mitte zwischen $s_{10}$
    und $s_{01}$, wie in Gleichung~\eqref{c02.eq.quadrat-mitte}. (b) Die Niveaulinien des
    Extremaleffekts $f_x^{+}$ aus Proposition~\ref{c02.prop.quadrat}: Er hat auf der
    rechten Kante den Wert $1$ und auf der linken den Wert $0$.}
\end{tpcfigure}

\begin{proposition}[Effekte des Quadrat-Bits]\label{c02.prop.quadrat}\claim[proof-included]{c02.cl.quadrat-effekte}{derived}
  Die Funktion $f(\vb{r}) = b + a_x r_x + a_y r_y$ ist genau dann ein Effekt auf
  $\StSp_{\mathrm{Q}}$, wenn
  \begin{equation}\label{c02.eq.quadrat-effekt}
    \abs{a_x} + \abs{a_y} \leq \min(b,\, 1-b) .
  \end{equation}
  Schreibt man einen Effekt als Punkt $(b, a_x, a_y)$, so ist
  $\EffAlg{\StSp_{\mathrm{Q}}}$ eine Doppelpyramide mit sechs Ecken. Diese sechs Ecken sind
  die Extremaleffekte:
  \begin{equation}\label{c02.eq.quadrat-extremal}
    \begin{aligned}
      0 &= (0,0,0), & f_x^{\pm} &= \bigl(\tfrac12, \pm\tfrac12, 0\bigr), \\
      \UnitEff{\StSp} &= (1,0,0), & f_y^{\pm} &= \bigl(\tfrac12, 0, \pm\tfrac12\bigr) .
    \end{aligned}
  \end{equation}
\end{proposition}
\begin{proof}
  Den kleinsten Wert nimmt $f$ in der Ecke an, in der $r_x$ das entgegengesetzte
  Vorzeichen von $a_x$ und $r_y$ das entgegengesetzte Vorzeichen von $a_y$ hat; er ist
  $b - \abs{a_x} - \abs{a_y}$. Der größte Wert ist entsprechend
  $b + \abs{a_x} + \abs{a_y}$. Beide liegen genau dann in $[0,1]$, wenn
  \eqref{c02.eq.quadrat-effekt} gilt.

  Jeder Effekt ist eine Mischung der sechs Punkte in \eqref{c02.eq.quadrat-extremal}. Mit
  $q = \abs{a_x} + \abs{a_y}$ ist nämlich
  \begin{equation*}
    (b, a_x, a_y) = (b - q)\,\UnitEff{\StSp} + 2\abs{a_x}\, f_x^{\pm} + 2\abs{a_y}\, f_y^{\pm} + (1 - b - q)\cdot 0 ,
  \end{equation*}
  wobei bei $f_x^{\pm}$ das Vorzeichen von $a_x$ und bei $f_y^{\pm}$ das von $a_y$ zu nehmen
  ist; ist $a_x = 0$ oder $a_y = 0$, so hat das betreffende Glied das Gewicht $0$ und
  entfällt. Die vier Gewichte sind wegen \eqref{c02.eq.quadrat-effekt} nicht negativ, und ihre
  Summe ist $1$. Ein Extremaleffekt muss daher einer der sechs Punkte sein.

  Dass alle sechs extremal sind, sieht man an den Stellen, an denen sie die Werte $0$ und
  $1$ annehmen. Sei $f_x^{+} = \lambda g + (1-\lambda) h$ mit Effekten $g$, $h$ und
  $0 < \lambda < 1$. Der Effekt $f_x^{+}(\vb{r}) = \tfrac12(1 + r_x)$ hat in den Ecken
  $(1,1)$ und $(1,-1)$ den Wert $1$ und in $(-1,1)$ den Wert $0$. Wie im Beweis von
  Proposition~\ref{c02.prop.bit} müssen $g$ und $h$ dort dieselben Werte haben. Drei
  Gleichungen $b + a_x + a_y = 1$, $b + a_x - a_y = 1$, $b - a_x + a_y = 0$ haben nur die
  Lösung $(b, a_x, a_y) = (\tfrac12, \tfrac12, 0)$; also ist $g = h = f_x^{+}$. Für
  $f_x^{-}$, $f_y^{+}$, $f_y^{-}$ geht es genauso. Die Nullfunktion ist extremal, weil ein
  Mittel zweier Funktionen mit Werten in $[0,1]$ nur dann überall $0$ ist, wenn beide
  überall $0$ sind; für $\UnitEff{\StSp}$ gilt dasselbe mit dem Wert $1$.
\end{proof}

Auch hier haben die Extremaleffekte eine einfache Bedeutung. Der Effekt
$f_x^{+} = \tfrac12(1 + r_x)$ hat in der Ecke $s_{ij}$ den Wert $i$: Er fragt, ob der
Kasten bei Eingabe $0$ die Ausgabe $1$ liefert. Entsprechend hat
$f_y^{+} = \tfrac12(1 + r_y)$ in $s_{ij}$ den Wert $j$ und fragt nach der Ausgabe bei
Eingabe $1$. $f_x^{-}$ und $f_y^{-}$ sind die Komplemente, $f_x^{-} = \UnitEff{\StSp} - f_x^{+}$
und $f_y^{-} = \UnitEff{\StSp} - f_y^{+}$. Plávala gibt dieselbe Menge als konvexe Hülle
dieser sechs Effekte an; sie ist in Abbildung~\ref{c02.fig.effektmengen}(b) gezeichnet.

\subsection{Das Qubit}\label{c02.s4.3}

Die dritte Theorie ist das kleinste nichttriviale System der Quantentheorie. Seine Zustände sind
Matrizen: komplexe $2\times 2$-Matrizen $\rho$ mit drei Eigenschaften,
\begin{equation}\label{c02.eq.qubit}\claim{c02.cl.qubit-zustaende}{source-backed}
  \rho\ \text{\tpcifenglish{Hermitian}{hermitesch}}, \qquad \rho \geq 0, \qquad \tr(\rho) = 1 .
\end{equation}
Hermitesch heißt: Die Diagonaleinträge sind reell, und die beiden anderen Einträge sind
zueinander komplex konjugiert. Die Bedingung $\rho \geq 0$ heißt für eine hermitesche
$2\times 2$-Matrix: Ihre beiden Eigenwerte sind nicht negativ. Solche Matrizen heißen
Dichtematrizen \srcref{plavala-2103.07469v2}{„Example: quantum theory“}.
Um sie zu überblicken, benutzen wir die Einheitsmatrix und die drei Pauli-Matrizen
\srcref{tpc-script-v0.2}{Vorlesung 17}:
\begin{equation}\label{c02.eq.pauli}\claim{c02.cl.pauli}{source-backed}
  \Id = \begin{pmatrix} 1 & 0 \\ 0 & 1 \end{pmatrix}, \quad
  \pauli{x} = \begin{pmatrix} 0 & 1 \\ 1 & 0 \end{pmatrix}, \quad
  \pauli{y} = \begin{pmatrix} 0 & -i \\ i & 0 \end{pmatrix}, \quad
  \pauli{z} = \begin{pmatrix} 1 & 0 \\ 0 & -1 \end{pmatrix} .
\end{equation}

Die folgende Proposition sagt, dass diese Matrizen eine Kugel bilden (Abbildung~\ref{c02.fig.blochzustaende}).

\begin{tpcfigure}{c02.fig.blochzustaende}
  \caption{Die Zustände des Qubits als Punkte der Einheitskugel. Die Punkte $\vb{n}$, $-\vb{n}$, $\vb{m}$, $-\vb{m}$ der Oberfläche sind reine Zustände. Der Zustand $\vb{r} = \tfrac12\vb{n}$ im Inneren ist die Mischung $\tfrac34\vb{n} + \tfrac14(-\vb{n})$. Der Mittelpunkt ist die Mitte zwischen $\vb{n}$ und $-\vb{n}$ und ebenso die Mitte zwischen $\vb{m}$ und $-\vb{m}$.}
\end{tpcfigure}

\begin{proposition}[Die Zustandskugel des Qubits]\label{c02.prop.bloch}\claim[proof-included]{c02.cl.bloch}{derived}
  \begin{enumerate}
    \item Jede hermitesche $2\times 2$-Matrix mit Spur $1$ lässt sich mit genau einem
      Vektor $\vb{r} = (r_x, r_y, r_z) \in \mathbb{R}^3$ schreiben als
      \begin{equation}\label{c02.eq.bloch}
        \rho(\vb{r}) = \tfrac12\bigl(\Id + r_x \pauli{x} + r_y \pauli{y} + r_z \pauli{z}\bigr)
        = \frac12 \begin{pmatrix} 1 + r_z & r_x - i r_y \\ r_x + i r_y & 1 - r_z \end{pmatrix} .
      \end{equation}
    \item Die Eigenwerte von $\rho(\vb{r})$ sind $\tfrac12(1 + \norm{\vb{r}})$ und
      $\tfrac12(1 - \norm{\vb{r}})$. Also ist $\rho(\vb{r})$ genau dann ein Zustand, wenn
      $\norm{\vb{r}} \leq 1$ ist.
    \item Mischungen bleiben Mischungen:
      $\rho(\lambda \vb{r} + (1-\lambda)\vb{s}) = \lambda \rho(\vb{r}) + (1-\lambda)\rho(\vb{s})$.
    \item Der Zustand $\rho(\vb{r})$ ist genau dann rein, wenn $\norm{\vb{r}} = 1$ ist.
  \end{enumerate}
\end{proposition}
\begin{proof}
  (1) Eine hermitesche Matrix hat reelle Diagonaleinträge; ist ihre Spur $1$, so lassen
  sie sich als $\tfrac12(1 + r_z)$ und $\tfrac12(1 - r_z)$ schreiben, wobei $r_z$ ihre
  Differenz ist. Der Eintrag links unten ist eine komplexe Zahl, die wir
  $\tfrac12(r_x + i r_y)$ nennen; der Eintrag rechts oben ist dann die konjugierte Zahl
  $\tfrac12(r_x - i r_y)$.

  (2) Die Eigenwerte $\kappa$ sind die Nullstellen von
  \begin{equation*}
    \det\bigl(\kappa\, \Id - \rho(\vb{r})\bigr)
    = \bigl(\kappa - \tfrac12(1 + r_z)\bigr)\bigl(\kappa - \tfrac12(1 - r_z)\bigr) - \tfrac14\bigl(r_x^2 + r_y^2\bigr)
    = \bigl(\kappa - \tfrac12\bigr)^2 - \tfrac14 \norm{\vb{r}}^2 .
  \end{equation*}
  Beide sind genau dann nicht negativ, wenn $\norm{\vb{r}} \leq 1$ ist.

  (3) Die rechte Seite von \eqref{c02.eq.bloch} ist eine Konstante plus ein Ausdruck, der
  linear in $\vb{r}$ ist.

  (4) Sei $\norm{\vb{r}} = 1$ und $\vb{r} = \lambda \vb{s} + (1-\lambda)\vb{t}$ mit
  $\norm{\vb{s}} \leq 1$, $\norm{\vb{t}} \leq 1$ und $0 < \lambda < 1$. Dann ist
  $1 = \vb{r}\cdot\vb{r} = \lambda\, \vb{r}\cdot\vb{s} + (1-\lambda)\, \vb{r}\cdot\vb{t}$.
  Ein Skalarprodukt ist höchstens so groß wie das Produkt der Längen, hier also höchstens
  $1$, und es ist nur dann gleich $1$, wenn beide Vektoren gleich sind. Also ist
  $\vb{s} = \vb{t} = \vb{r}$, und $\rho(\vb{r})$ ist rein. Ist dagegen
  $q = \norm{\vb{r}} < 1$, so wählen wir einen Einheitsvektor $\vb{n}$ mit
  $\vb{r} = q\,\vb{n}$ und schreiben
  $\vb{r} = \tfrac12(1+q)\,\vb{n} + \tfrac12(1-q)\,(-\vb{n})$. Das ist eine Mischung zweier
  verschiedener Zustände mit Gewichten echt zwischen $0$ und $1$.
\end{proof}

Wegen (3) dürfen wir den Zustand $\rho(\vb{r})$ mit dem Punkt $\vb{r}$ gleichsetzen. Der
Zustandsraum des Qubits ist dann die Einheitskugel
$\StSp_{\mathrm{B}} = \{ \vb{r} \in \mathbb{R}^3 : \norm{\vb{r}} \leq 1 \}$. Wir nennen
sie die Zustandskugel des Qubits; ihre reinen Zustände sind die Punkte der Oberfläche
(Abbildung~\ref{c02.fig.zustandsraeume}(c)). Auch hier ist die Zerlegung in reine Zustände
nicht eindeutig: Der Mittelpunkt ist die Mitte zwischen $\vb{n}$ und $-\vb{n}$ für jeden
Einheitsvektor $\vb{n}$.

Die Effekte gewinnt man wie beim Quadrat, nur mit einer anderen Länge für $\vb{a}$. Auch
sie lassen sich als Matrizen schreiben. Abbildung~\ref{c02.fig.blochkugel} zeigt vorab,
wie ein Extremaleffekt der Kugel aussieht.

\begin{tpcfigure}{c02.fig.blochkugel}
  \caption{Die Zustandskugel des Qubits mit dem Extremaleffekt
    $f_{\vb{n}}(\vb{r}) = \tfrac12(1 + \vb{n}\cdot\vb{r})$ aus
    Proposition~\ref{c02.prop.qubit-effekte} für $\vb{n} = (0,0,1)$, also
    $f_{\vb{n}}(\vb{r}) = \tfrac12(1 + r_z)$. Die Punkte gleicher Wahrscheinlichkeit
    bilden waagerechte Ebenen. Die Ebene $f = 1$ berührt die Kugel im reinen Zustand
    $\vb{n}$ (oben), die Ebene $f = 0$ im reinen Zustand $-\vb{n}$ (unten). Der Zustand $x$
    liegt auf der Höhe $r_z = \tfrac12$; dort ist $f = \tfrac34$. Bei einem nicht
    konstanten Effekt, der nicht extremal ist, berührt höchstens eine der beiden Ebenen
    $f = 0$ und $f = 1$ die Kugel; eine von ihnen kann sie noch berühren. Ein konstanter
    Effekt hat überall denselben Wert; seine Niveaumengen sind keine Ebenen.}
\end{tpcfigure}

\begin{proposition}[Effekte des Qubits]\label{c02.prop.qubit-effekte}\claim[proof-included]{c02.cl.qubit-effekte}{derived}
  \begin{enumerate}
    \item Die Funktion $f(\vb{r}) = b + \vb{a}\cdot\vb{r}$ ist genau dann ein Effekt auf
      $\StSp_{\mathrm{B}}$, wenn
      \begin{equation}\label{c02.eq.qubit-effekt}
        \norm{\vb{a}} \leq \min(b,\, 1-b) .
      \end{equation}
    \item Mit der hermiteschen Matrix
      $\op{\EffOp} = b\, \Id + a_x \pauli{x} + a_y \pauli{y} + a_z \pauli{z}$ ist
      \begin{equation}\label{c02.eq.spurformel}
        f(\vb{r}) = b + \vb{a}\cdot\vb{r} = \tr\bigl(\rho(\vb{r})\, \op{\EffOp}\bigr) .
      \end{equation}
      Die Eigenwerte von $\op{\EffOp}$ sind $b + \norm{\vb{a}}$ und $b - \norm{\vb{a}}$;
      Bedingung~\eqref{c02.eq.qubit-effekt} sagt, dass beide in $[0,1]$ liegen.
    \item Die Extremaleffekte sind $0$, $\UnitEff{\StSp}$ und, für jeden Einheitsvektor
      $\vb{n}$, der Effekt
      \begin{equation}\label{c02.eq.qubit-extremal}
        f_{\vb{n}}(\vb{r}) = \tfrac12\bigl(1 + \vb{n}\cdot\vb{r}\bigr) .
      \end{equation}
  \end{enumerate}
\end{proposition}
\begin{proof}
  (1) Auf der Kugel ist $\vb{a}\cdot\vb{r}$ höchstens $\norm{\vb{a}}$ und mindestens
  $-\norm{\vb{a}}$; beide Werte werden angenommen, für $\vb{a} \neq 0$ in den Punkten
  $\pm\vb{a}/\norm{\vb{a}}$. Der kleinste Wert von $f$ ist also $b - \norm{\vb{a}}$, der
  größte $b + \norm{\vb{a}}$.

  (2) Durch Ausmultiplizieren der Matrizen in \eqref{c02.eq.pauli} findet man: Jede
  Pauli-Matrix hat Spur $0$, ihr Quadrat ist $\Id$, und das Produkt zweier verschiedener
  Pauli-Matrizen hat Spur $0$. Also ist $\tr(\Id) = 2$, die Spur des Quadrats einer
  Pauli-Matrix ist $2$, und im Produkt $\rho(\vb{r})\,\op{\EffOp}$ tragen nur die Glieder mit
  $\Id\cdot\Id$ und mit zwei gleichen Pauli-Matrizen zur Spur bei:
  $\tr(\rho(\vb{r})\,\op{\EffOp}) = \tfrac12\bigl(2b + 2 a_x r_x + 2 a_y r_y + 2 a_z r_z\bigr)$.
  Die Eigenwerte von $\op{\EffOp}$ erhält man wie in Proposition~\ref{c02.prop.bloch}(2).

  (3) Sei $q = \norm{\vb{a}}$. Für $q = 0$ ist
  $(b, 0) = b\, \UnitEff{\StSp} + (1 - b)\cdot 0$, eine Mischung von $\UnitEff{\StSp}$ und
  der Nullfunktion. Für $q > 0$ setzen wir $\vb{n} = \vb{a}/q$; dann ist
  \begin{equation*}
    (b, \vb{a}) = (b - q)\, \UnitEff{\StSp} + 2q\, \bigl(\tfrac12, \tfrac12\vb{n}\bigr) + (1 - b - q)\cdot 0 ,
  \end{equation*}
  eine Mischung mit nicht negativen Gewichten der Summe $1$. Extremaleffekte kann es also
  nur unter den genannten geben. Der Effekt $f_{\vb{n}}$ hat im Punkt $\vb{n}$ den Wert
  $1$ und im Punkt $-\vb{n}$ den Wert $0$. Ist $f_{\vb{n}} = \lambda g + (1-\lambda) h$ mit
  Effekten $g$, $h$ und $0 < \lambda < 1$, so haben $g$ und $h$ dort dieselben Werte. Für
  $g = (b, \vb{a})$ heißt das $b + \vb{a}\cdot\vb{n} = 1$ und $b - \vb{a}\cdot\vb{n} = 0$,
  also $b = \tfrac12$ und $\vb{a}\cdot\vb{n} = \tfrac12$. Wegen
  \eqref{c02.eq.qubit-effekt} ist $\norm{\vb{a}} \leq \tfrac12$, und das lässt nur
  $\vb{a} = \tfrac12\vb{n}$ zu. Also ist $g = h = f_{\vb{n}}$. Für $0$ und
  $\UnitEff{\StSp}$ gilt das Argument aus Proposition~\ref{c02.prop.quadrat}.
\end{proof}

Der Extremaleffekt $f_{\vb{n}}$ fragt: „Liegt der reine Zustand $\vb{n}$ vor?“ Im Zustand
$\vb{n}$ ist die Antwort sicher Ja, im gegenüberliegenden Zustand $-\vb{n}$ sicher Nein, und
in jedem anderen Zustand ist sie zufällig. Seine Matrix ist
$\tfrac12(\Id + n_x\pauli{x} + n_y\pauli{y} + n_z\pauli{z}) = \rho(\vb{n})$, die Dichtematrix
dieses reinen Zustands. In Abbildung~\ref{c02.fig.blochkugel} ist er als Schar paralleler
Ebenen gezeichnet.

\subsection{Was die drei unterscheidet}\label{c02.s4.4}

In Koordinaten mit dem Nullpunkt in der Mitte sehen die drei Theorien fast gleich aus. Ein
Effekt ist jedes Mal eine Funktion $b + \vb{a}\cdot\vb{r}$ mit der Bedingung, dass eine
Länge von $\vb{a}$ höchstens $\min(b, 1-b)$ ist. Nur die Länge wird verschieden gemessen:
beim Bit als Betrag $\abs{a}$ (Gleichung~\eqref{c02.eq.bit-zentriert}), beim Quadrat-Bit
als Summe $\abs{a_x} + \abs{a_y}$ (Gleichung~\eqref{c02.eq.quadrat-effekt}), beim Qubit als
gewöhnliche Länge $\norm{\vb{a}}$ (Gleichung~\eqref{c02.eq.qubit-effekt}).
Tabelle~\ref{c02.tab.vergleich} stellt zusammen, was daraus folgt.

\begin{tpctable}{c02.tab.vergleich}
  \caption{Die drei Theorien im Vergleich. Die Einträge sind die Ergebnisse der
    Propositionen~\ref{c02.prop.bit}, \ref{c02.prop.quadrat}, \ref{c02.prop.bloch} und
    \ref{c02.prop.qubit-effekte} sowie der Gleichungen~\eqref{c02.eq.bit} und
    \eqref{c02.eq.quadrat-mitte}.}
  \begin{tabular}{llll}
    \toprule
      & Bit & Quadrat-Bit & Qubit \\
    \midrule
    Zustandsraum & Strecke & Quadrat & Kugel \\
    reine Zustände & zwei & vier & Oberfläche \\
    Extremaleffekte & vier & sechs & $0$, $\UnitEff{\StSp}$, alle $f_{\vb{n}}$ \\
    Zerlegung eindeutig & ja & nein & nein \\
    klassisch (\defref{D11}) & ja & nein & nein \\
    \bottomrule
  \end{tabular}
\end{tpctable}

Dass das Qubit keine klassische Theorie ist, folgt wie beim Quadrat aus der mehrdeutigen
Zerlegung. Worin sich Quadrat-Bit und Qubit voneinander unterscheiden, zeigt die Aufgabe in
Abschnitt~\ref{c02.s6}.

\section{Messungen und die Hypothese der Nichteinschränkung}\label{c02.s5}

Eine Ja-Nein-Frage hat zwei mögliche Ergebnisse. Für Messungen mit mehr Ergebnissen braucht
man einen weiteren Begriff.

\tpcdef{D8}

\tpcdef{D9}

Kanäle beschreiben allgemein, wie aus Zuständen eines Systems Zustände eines anderen
werden. Dieses Kapitel braucht sie nur für \defref{D9}, und davon nur die zweite Hälfte:
Eine Messung mit $n$ Ergebnissen ist eine Liste von Effekten $f_1, \dots, f_n$, deren Summe
der Einseffekt ist. Im Zustand $x$ tritt das $i$-te Ergebnis mit der Wahrscheinlichkeit
$f_i(x)$ ein, und diese Zahlen bilden eine Wahrscheinlichkeitsverteilung:
\begin{equation}\label{c02.eq.messung}\claim{c02.cl.messung}{derived}
  f_i(x) \geq 0 \quad (i = 1, \dots, n),
  \qquad f_1(x) + \dots + f_n(x) = \UnitEff{\StSp}(x) = 1 .
\end{equation}
Für $n = 2$ ist das die Ja-Nein-Frage mit $f_1 = f$ und $f_2 = \UnitEff{\StSp} - f$ aus
Gleichung~\eqref{c02.eq.komplement}
\srcref{plavala-2103.07469v2}{„Measurements“}.

\begin{tpcfigure}{c02.fig.messung}
  \caption{Eine Messung mit drei Ergebnissen auf dem klassischen Bit. Die Effekte sind $f_1 = (\tfrac12, 0)$, $f_2 = (\tfrac12, \tfrac12)$ und $f_3 = (0, \tfrac12)$ in den Koordinaten $(q_0, q_1)$; ihre Summe ist $(1,1) = \UnitEff{\StSp}$. Über jedem Zustand sind die drei Wahrscheinlichkeiten übereinander aufgetragen; zusammen füllen sie die Höhe $1$. Im Zustand $x = (\tfrac34, \tfrac14)$ sind es $\tfrac38$, $\tfrac12$ und $\tfrac18$.}
\end{tpcfigure}

Bisher war ein Effekt eine Funktion mit bestimmten Eigenschaften. Ob sich zu jeder solchen
Funktion auch ein Experiment angeben lässt, ist eine eigene Frage.

\tpcdef{D12}

% GAP-TJO: Ob die Hypothese der Nichteinschränkung im Kurs Teil von Postulat P3 ist oder
% eine eigene Annahme, sagt docs/plan/BRIEF.md nicht (BRIEF.md:53).  Der Text unten nennt
% sie nur als Hypothese dieses Kapitels.
Dieses Kapitel setzt die Hypothese überall voraus: In Abschnitt~\ref{c02.s6} wird über
alle Effekte aus $\EffAlg{\StSp}$ optimiert. Lässt eine Theorie nur einen Teil der
Effekte zu, so ist das dort berechnete Optimum eine obere Schranke und muss nicht
erreichbar sein. Die Hypothese spricht außerdem nur von Ja-Nein-Fragen. Dass sich auch jede
Liste von Effekten mit Summe $\UnitEff{\StSp}$ als Messung ausführen lässt, wäre eine
weitere Annahme; die Unterscheidungsaufgabe in Abschnitt~\ref{c02.s6} benutzt nur
Messungen mit zwei Ergebnissen und braucht sie nicht\zoomref{c02.v2}. Die Messung mit drei
Ergebnissen in Abbildung~\ref{c02.fig.messung} ist ein Beispiel zu \defref{D9} und kommt
in der Aufgabe nicht vor.

\begin{vertiefung}[id=c02.v2,
    titel={Zwei Kegel und die Ebene auf Höhe eins},
    wozu={Zustände und Effekte leben in zwei verschiedenen Vektorräumen, die zueinander
      dual sind. Hier werden beide Räume und ihre Kegel eingeführt und für die drei
      Theorien ausgerechnet.},
    voraussetzung={Kern bis Abschnitt~\ref{c02.s5}; Vektorräume, lineare Abbildungen. Der
      Kern braucht diesen Block nicht; die Vertiefungen~\ref{c02.v3} und \ref{c02.v5} bauen
      auf ihm auf.},
    zurueck={Ein Zustand ist ein Punkt eines Kegels auf Höhe eins; ein Effekt ist ein
      Punkt eines anderen Kegels zwischen $0$ und dem Einseffekt.}]
  \tpcdef{D57}

  \tpcdef{D59}

  Der Kegel $\AffPos{\StSp}$ besteht aus Funktionen auf Zuständen; in ihm liegen die
  Effekte. Der Kegel $\AffDualPos{\StSp}$ besteht aus Funktionalen auf diesen Funktionen;
  in ihm liegen die Zustände. Die beiden sind nicht zu verwechseln. Wo die Zustände in
  ihrem Kegel liegen, sagt ein Satz von Plávala
  \srcref{plavala-2103.07469v2}{„Duality between the state space and the effect algebra“}:

  \begin{theorem}[Zustände auf Höhe eins]\label{c02.thm.hoehe}\claim[proof-omitted]{c02.cl.hoehe}{source-quoted}
    Für jeden Zustandsraum $\StSp$ gilt, bis auf einen Isomorphismus,
    \begin{equation}\label{c02.eq.hoehe}
      \StSp = \bigl\{ \varphi \in \AffDualPos{\StSp} : \pair{\varphi}{\UnitEff{\StSp}} = 1 \bigr\} .
    \end{equation}
  \end{theorem}

  Der Isomorphismus ist die Auswertung aus \defref{D59}: Der Zustand $x$ wird zum
  Funktional $f \mapsto f(x)$. Die Bedingung $\pair{\varphi}{\UnitEff{\StSp}} = 1$ ist die
  Normierung; sie schneidet aus dem Kegel eine Ebene aus, die „Höhe eins“.

  \begin{lemma}[Der Kegel der unnormierten Zustände]\label{c02.lem.basis}\claim[proof-included]{c02.cl.basis}{derived}
    Jedes $\psi \in \AffDualPos{\StSp}$ mit $\psi \neq 0$ lässt sich auf genau eine Weise
    als $\psi = t\, x$ mit einer Zahl $t > 0$ und einem Zustand $x \in \StSp$ schreiben.
    Dabei ist $t = \pair{\psi}{\UnitEff{\StSp}}$.
  \end{lemma}
  \begin{proof}
    Zu jedem $f \in \Aff{\StSp}$ gibt es eine Zahl $M \geq 0$ mit
    $-M \leq f(x) \leq M$ für alle $x \in \StSp$, denn $f$ ist nach
    Lemma~\ref{c02.lem.koordinaten} stetig und $\StSp$ ist abgeschlossen und beschränkt.
    Die Funktionen $M\,\UnitEff{\StSp} - f$ und $M\,\UnitEff{\StSp} + f$ liegen dann in
    $\AffPos{\StSp}$, und für $\psi \in \AffDualPos{\StSp}$ folgt
    $-M\, \pair{\psi}{\UnitEff{\StSp}} \leq \pair{\psi}{f} \leq M\, \pair{\psi}{\UnitEff{\StSp}}$.
    Wäre $t = \pair{\psi}{\UnitEff{\StSp}} = 0$, so wäre also $\pair{\psi}{f} = 0$ für
    alle $f$ und damit $\psi = 0$. Für $\psi \neq 0$ ist deshalb $t > 0$, das Funktional
    $\psi / t$ ist positiv und normiert, und nach Theorem~\ref{c02.thm.hoehe} ist es ein
    Zustand $x$. Ist $\psi = t' x'$ eine zweite solche Darstellung, so liefert die Paarung
    mit $\UnitEff{\StSp}$ zuerst $t' = t$ und dann $x' = x$.
  \end{proof}

  $\AffDualPos{\StSp}$ ist also der Kegel der Strahlen vom Nullpunkt durch $\StSp$: die
  Menge der unnormierten Zustände. Den zweiten Kegel beschreibt Plávala so:

  \begin{lemma}[Positive Funktionen und positive Funktionale]\label{c02.lem.dualkegel}\claim[proof-omitted]{c02.cl.dualkegel}{source-quoted}
    Sei $f \in \Aff{\StSp}$. Dann ist $f \in \AffPos{\StSp}$ genau dann, wenn
    $\pair{\psi}{f} \geq 0$ für alle $\psi \in \AffDualPos{\StSp}$ gilt.
  \end{lemma}

  Der Kegel der positiven affinen Funktionen ist demnach der duale Kegel zum Kegel der
  unnormierten Zustände, und die Effekte sind darin die Funktionen zwischen $0$ und
  $\UnitEff{\StSp}$. An dieser Stelle steht in der Rahmung des Kurses der Satz
  % GAP-TJO: Wie TJO „elements of the dual space to states“ (BRIEF.md:46) gelesen haben
  % will, sagt der Brief nicht (evidence/gaps.json, G-0010).  Der Satz wird deshalb nur
  % zitiert; die mathematische Aussage danach ist die des Kapitels und keine Auslegung
  % seines Satzes.  Die deutsche Fassung des Satzes ist eine Übersetzung.
  \begin{equation}\label{c02.eq.rahmen-dual}\claim{c02.cl.rahmen-dual}{author-stipulated}
    \begin{aligned}
      &\text{\tpcifenglish{Measurements/instruments/POVMs are then}{Messungen/Instrumente/POVMs sind dann}} \\
      &\text{\tpcifenglish{elements of the dual space to states.}{Elemente des Dualraums zu den Zuständen.}}
    \end{aligned}
  \end{equation}
  Dieser Satz wird hier zitiert und nicht ausgelegt. Getrennt davon hält dieser Block als
  mathematische Aussage fest: Ein einzelner Effekt ist ein Element von
  $\Aff{\StSp}$ mit $0 \leq f \leq \UnitEff{\StSp}$, und $\Aff{\StSp}$ und der von den
  Zuständen aufgespannte Raum $\AffDual{\StSp}$ sind zueinander dual. Eine Messung ist eine
  Liste solcher Effekte (\defref{D9}).

  \emph{In Koordinaten.} Für $\StSp \subset \mathbb{R}^d$ wie in
  Lemma~\ref{c02.lem.koordinaten} schreiben wir den Zustand $\vb{r}$ als Vektor
  $(1, \vb{r}) \in \mathbb{R}^{d+1}$ und die Funktion $b + \vb{a}\cdot\vb{r}$ als Vektor
  $(b, \vb{a})$. Die Paarung ist dann das gewöhnliche Skalarprodukt dieser beiden Vektoren,
  $b\cdot 1 + \vb{a}\cdot\vb{r}$. Auf den Vektoren $(1, \vb{r})$ sind Effekte linear; auf
  den Punkten $\vb{r}$ selbst sind sie nur affin. Die beiden Kegel der drei Theorien sind
  \begin{equation}\label{c02.eq.kegel-bit}\claim{c02.cl.kegel-bit}{derived}
    \begin{aligned}
      \AffDualPos{\Simplex{2}} &= \{ (t, R) : t \geq \abs{R} \}, \\
      \AffPos{\Simplex{2}} &= \{ (b, a) : b \geq \abs{a} \},
    \end{aligned}
  \end{equation}
  \begin{equation}\label{c02.eq.kegel-quadrat}\claim{c02.cl.kegel-quadrat}{derived}
    \begin{aligned}
      \AffDualPos{\StSp_{\mathrm{Q}}} &= \{ (t, R_x, R_y) : t \geq \max(\abs{R_x}, \abs{R_y}) \}, \\
      \AffPos{\StSp_{\mathrm{Q}}} &= \{ (b, a_x, a_y) : b \geq \abs{a_x} + \abs{a_y} \},
    \end{aligned}
  \end{equation}
  \begin{equation}\label{c02.eq.kegel-qubit}\claim{c02.cl.kegel-qubit}{derived}
    \begin{aligned}
      \AffDualPos{\StSp_{\mathrm{B}}} &= \{ (t, \vb{R}) : t \geq \norm{\vb{R}} \}, \\
      \AffPos{\StSp_{\mathrm{B}}} &= \{ (b, \vb{a}) : b \geq \norm{\vb{a}} \} ,
    \end{aligned}
  \end{equation}
  das Bit dabei in der Koordinate $r$ aus Gleichung~\eqref{c02.eq.bit-zentriert}.
  In der ersten Zeile steht jeweils die Menge der Vektoren $t\,(1, \vb{r})$ mit $t \geq 0$ und
  $\vb{r} \in \StSp$, in der zweiten die Menge der Funktionen, deren kleinster Wert auf $\StSp$
  nicht negativ ist; die kleinsten Werte wurden in Abschnitt~\ref{c02.s4} berechnet. Beim
  Bit und beim Qubit sind die beiden Mengen in diesen Koordinaten gleich. Beim Quadrat-Bit
  sind sie verschieden: Der Vektor $(1, 1, 1)$ liegt in der ersten Menge und nicht in der
  zweiten (Abbildung~\ref{c02.fig.kegel}). Eine Aussage über andere Koordinaten oder andere
  Skalarprodukte ist das nicht.

  \begin{tpcfigure}{c02.fig.kegel}
    \caption{Die beiden Kegel des Quadrat-Bits aus \eqref{c02.eq.kegel-quadrat}. (a) Der
      Kegel $\AffDualPos{\StSp}$ der unnormierten Zustände. Der Zustandsraum $\StSp$ ist
      sein Schnitt mit der Ebene $t = 1$; jeder Strahl vom Nullpunkt durch den Kegel trifft
      $\StSp$ in genau einem Zustand $x$. (b) Der Kegel $\AffPos{\StSp}$ der positiven
      affinen Funktionen; nach oben ist $b$ aufgetragen. Die Effektmenge ist die
      gestrichelte Doppelpyramide zwischen $0$ und $\UnitEff{\StSp}$ aus
      Abbildung~\ref{c02.fig.effektmengen}(b). Der Kegel in (a) hat einen quadratischen
      Querschnitt mit Seiten parallel zu den Achsen, der in (b) einen um ein Achtel einer
      vollen Drehung gedrehten und schmaleren: Die beiden Kegel sind nicht gleich.}
  \end{tpcfigure}

  Der Einseffekt spielt in $\Aff{\StSp}$ eine besondere Rolle: Jede affine Funktion liegt
  unter einem Vielfachen von ihm, wie im Beweis von Lemma~\ref{c02.lem.basis} benutzt. Ein
  solches Element heißt Ordnungseinheit\zoomref{c02.f1}.
\end{vertiefung}

\section{Zwei Präparationen unterscheiden}\label{c02.s6}

Jetzt zur Aufgabe, an der sich die drei Theorien messen lassen. Gegeben sind zwei Zustände
$x_1$ und $x_2$ desselben Zustandsraums $\StSp$ und zwei Zahlen $\lambda_1, \lambda_2 \geq 0$
mit $\lambda_1 + \lambda_2 = 1$. Jemand präpariert mit Wahrscheinlichkeit $\lambda_1$ den
Zustand $x_1$ und mit Wahrscheinlichkeit $\lambda_2$ den Zustand $x_2$ und sagt nicht,
welchen. Wir dürfen eine Ja-Nein-Frage $f$ stellen und sollen dann raten. Die Verabredung
ist: Bei Ja raten wir $x_1$, bei Nein $x_2$.

Richtig liegen wir in zwei Fällen: Es wurde $x_1$ präpariert und die Antwort ist Ja, oder
es wurde $x_2$ präpariert und die Antwort ist Nein. Die Erfolgswahrscheinlichkeit ist
deshalb
\begin{equation}\label{c02.eq.erfolg}\claim{c02.cl.erfolg}{source-backed}
  P(f) = \lambda_1\, f(x_1) + \lambda_2\, \bigl(1 - f(x_2)\bigr) .
\end{equation}
Gesucht ist der größte Wert $P_{\mathrm{opt}}$, den $P(f)$ annimmt, wenn $f$ alle Effekte
durchläuft \srcref{plavala-2103.07469v2}{„Order unit norm, base norm and discrimination tasks“}.

\begin{tpcfigure}{c02.fig.aufgabe}
  \caption{Die Aufgabe. Links: Präpariert wird $x_1$ oder $x_2$, dann wird die Frage $f$ gestellt und nach der Antwort geraten. Rechts: die Wahrscheinlichkeiten der vier möglichen Fälle, in der oberen Zeile für die Antwort Ja, in der unteren für Nein. In den beiden stark umrandeten Fällen ist der Tipp richtig; ihre Summe ist $P(f)$ aus Gleichung~\eqref{c02.eq.erfolg}.}
\end{tpcfigure}

Zwei Effekte kommen ohne Messung aus. Die Nullfunktion antwortet immer Nein; wir raten dann
immer $x_2$ und liegen mit Wahrscheinlichkeit $\lambda_2$ richtig. Der Einseffekt antwortet
immer Ja; wir raten immer $x_1$ und liegen mit Wahrscheinlichkeit $\lambda_1$ richtig.
Deshalb gilt in jeder Theorie
\begin{equation}\label{c02.eq.schranken}\claim{c02.cl.schranken}{derived}
  \max(\lambda_1, \lambda_2) \leq P_{\mathrm{opt}} \leq 1 .
\end{equation}

\begin{proposition}[Wann gelingt die Unterscheidung sicher?]\label{c02.prop.perfekt}\claim[proof-included]{c02.cl.perfekt}{derived}
  Seien $\lambda_1 > 0$ und $\lambda_2 > 0$. Dann ist $P(f) = 1$ genau dann, wenn
  $f(x_1) = 1$ und $f(x_2) = 0$ ist.
\end{proposition}
\begin{proof}
  Aus \eqref{c02.eq.erfolg} und $\lambda_1 + \lambda_2 = 1$ folgt
  $1 - P(f) = \lambda_1\bigl(1 - f(x_1)\bigr) + \lambda_2\, f(x_2)$. Beide Summanden sind
  nicht negativ; die Summe ist nur dann $0$, wenn beide $0$ sind.
\end{proof}

Zwei Zustände heißen in diesem Kapitel \emph{sicher unterscheidbar}, wenn es einen Effekt
mit $f(x_1) = 1$ und $f(x_2) = 0$ gibt. Auf die Vorgaben $\lambda_1$, $\lambda_2$ kommt es
dabei nicht an, solange beide positiv sind. Ist eine von ihnen $0$, so erreicht schon
Raten den Wert $P = 1$, und der Wert sagt nichts über die beiden Zustände.

Die Funktion $P$ ist auf der Menge der Effekte affin: In den Koordinaten $(b, \vb{a})$ aus
Lemma~\ref{c02.lem.koordinaten} ist
\begin{equation}\label{c02.eq.erfolg-koordinaten}\claim{c02.cl.erfolg-koordinaten}{derived}
  P(f) = \lambda_2 + \delta\, b + \vb{v}\cdot\vb{a},
  \qquad \delta = \lambda_1 - \lambda_2, \quad \vb{v} = \lambda_1 \vb{r}_1 - \lambda_2 \vb{r}_2 ,
\end{equation}
wobei $\vb{r}_1$ und $\vb{r}_2$ die Koordinaten von $x_1$ und $x_2$ sind. Sind $\delta$
und $\vb{v}$ nicht beide $0$, so ist $P$ nicht konstant; die Gleichung
$\delta\, b + \vb{v}\cdot\vb{a} = c$ beschreibt dann für jede Zahl $c$ eine Hyperebene im
Koordinatenraum $\mathbb{R}^{d+1}$ der Effekte, diese Hyperebenen sind parallel, und die
Niveaumengen von $P$ sind ihre Schnitte mit der Effektmenge. In den Koordinaten $(b, a)$
des Bits sind das Geraden, beim Quadrat-Bit Ebenen. Sind $\delta = 0$ und $\vb{v} = 0$, so
ist $P$ konstant, und die ganze Effektmenge ist eine einzige Niveaumenge. Eine affine
Funktion nimmt ihren größten Wert am Rand an, und zwar in einer Ecke
(Abbildung~\ref{c02.fig.optimum})\zoomref{c02.v3}:

\begin{theorem}[Ein Extremaleffekt erreicht das Optimum]\label{c02.thm.extremal}\claim[proof-linked]{c02.cl.extremal}{derived}
  Sei $\StSp$ ein Zustandsraum, und seien $x_1, x_2 \in \StSp$ und
  $\lambda_1, \lambda_2 \geq 0$ mit $\lambda_1 + \lambda_2 = 1$ gegeben. Dann gibt es einen
  Extremaleffekt $f$ mit $P(f) = P_{\mathrm{opt}}$.
\end{theorem}

Der Satz sagt nicht, dass jeder optimale Effekt extremal ist. In
Abbildung~\ref{c02.fig.optimum}(b) sind zwei Extremaleffekte optimal und mit ihnen alle
ihre Mischungen.

\begin{tpcfigure}{c02.fig.optimum}
  \caption{Die Erfolgswahrscheinlichkeit $P$ als Funktion auf der Effektmenge; die dünnen
    Linien sind Niveaulinien von $P$. (a) Klassisches Bit mit $x_1 = (\tfrac34, \tfrac14)$,
    $x_2 = (\tfrac14, \tfrac34)$ und $\lambda_1 = \lambda_2 = \tfrac12$: Hier ist
    $P = \tfrac12 + \tfrac14(q_0 - q_1)$, und der größte Wert $\tfrac34$ wird nur in der
    Ecke $(1,0)$ angenommen. (b) Quadrat-Bit mit $x_1 = (1,1)$, $x_2 = (-1,-1)$ und
    $\lambda_1 = \lambda_2 = \tfrac12$: Hier ist $P = \tfrac12 + a_x + a_y$ unabhängig von
    $b$; gezeichnet ist der Schnitt $b = \tfrac12$ durch die Doppelpyramide. Der größte
    Wert $1$ wird auf der ganzen Kante zwischen $f_x^{+}$ und $f_y^{+}$ angenommen. Optimale
    Extremaleffekte sind eingekreist.}
\end{tpcfigure}

\begin{vertiefung}[id=c02.v3,
    titel={Warum ein Extremaleffekt genügt},
    wozu={Der Kern benutzt Theorem~\ref{c02.thm.extremal} ohne Beweis. Hier wird er für
      jeden Zustandsraum geführt.},
    voraussetzung={Vertiefungen~\ref{c02.v1} und \ref{c02.v2}; stetige Funktionen auf
      abgeschlossenen, beschränkten Mengen. Der Kern rechnet das Optimum für die drei
      Theorien in Proposition~\ref{c02.prop.optimum} direkt aus und braucht diesen Block
      nicht.},
    zurueck={Um zwei Präparationen bestmöglich zu unterscheiden, genügt es, die
      Extremaleffekte durchzugehen.}]
  \begin{lemma}[Die Effektmenge ist abgeschlossen, beschränkt und konvex]\label{c02.lem.kompakt}\claim[proof-included]{c02.cl.kompakt}{derived}
    Sei $\StSp \subset \mathbb{R}^d$ wie in Lemma~\ref{c02.lem.koordinaten}. In den
    Koordinaten $(b, \vb{a}) \in \mathbb{R}^{d+1}$ ist $\EffAlg{\StSp}$ eine
    abgeschlossene, beschränkte und konvexe Menge. Sie erfüllt damit selbst die
    Bedingungen (S1) bis (S5) aus \defref{D3}.
  \end{lemma}
  \begin{proof}
    $\EffAlg{\StSp}$ ist die Menge aller $(b, \vb{a})$, die für jedes $\vb{r} \in \StSp$
    die beiden Ungleichungen $0 \leq b + \vb{a}\cdot\vb{r} \leq 1$ erfüllen. Jede einzelne
    Ungleichung beschreibt eine abgeschlossene, konvexe Menge, und ein Durchschnitt solcher
    Mengen ist wieder von dieser Art. Beschränkt ist $\EffAlg{\StSp}$, weil $b$ und
    $\vb{a}$ nach dem Beweis von Lemma~\ref{c02.lem.koordinaten} die Lösung eines linearen
    Gleichungssystems sind, dessen rechte Seiten $f(\vb{v}_i)$ in $[0,1]$ liegen.
  \end{proof}

  \begin{proof}[Beweis von Theorem~\ref{c02.thm.extremal}]
    Nach Gleichung~\eqref{c02.eq.erfolg-koordinaten} ist $P$ eine stetige Funktion der
    Koordinaten $(b, \vb{a})$. Auf der abgeschlossenen und beschränkten Menge
    $\EffAlg{\StSp}$ nimmt sie ihren größten Wert in einem Effekt $f^{*}$ an. Nach
    Lemma~\ref{c02.lem.kompakt} gilt Theorem~\ref{c02.thm.zerlegung} auch für
    $\EffAlg{\StSp}$: Es gibt Extremaleffekte $g_1, \dots, g_k$ und Gewichte
    $\kappa_i \geq 0$ mit Summe $1$ und $f^{*} = \kappa_1 g_1 + \dots + \kappa_k g_k$. Da $P$ affin
    ist, folgt wie in Gleichung~\eqref{c02.eq.mehrfach}
    \begin{equation*}
      P_{\mathrm{opt}} = P(f^{*}) = \kappa_1 P(g_1) + \dots + \kappa_k P(g_k) .
    \end{equation*}
    Kein $P(g_i)$ ist größer als $P_{\mathrm{opt}}$. Ein gewichtetes Mittel solcher Zahlen
    ist nur dann gleich $P_{\mathrm{opt}}$, wenn jedes $g_i$ mit $\kappa_i > 0$ den Wert
    $P(g_i) = P_{\mathrm{opt}}$ hat. Mindestens ein solches $g_i$ gibt es.
  \end{proof}

  Liegt $\StSp$ in einem affinen Teilraum kleinerer Dimension, so wendet man das Argument
  in Koordinaten dieses Teilraums an.
\end{vertiefung}

Für die drei Theorien lässt sich das Optimum ausrechnen.

\begin{proposition}[Das Optimum in den drei Theorien]\label{c02.prop.optimum}\claim[proof-included]{c02.cl.optimum}{derived}
  Mit $\delta = \lambda_1 - \lambda_2$ und $\vb{v} = \lambda_1 \vb{r}_1 - \lambda_2 \vb{r}_2$
  wie in Gleichung~\eqref{c02.eq.erfolg-koordinaten} gilt:
  \begin{align}
    \text{\tpcifenglish{classical bit}{klassisches Bit}}\ (r = \pw_0 - \pw_1):\quad
      P_{\mathrm{opt}} &= \tfrac12\bigl(1 + \max(\abs{\delta},\, \abs{v})\bigr),
      \label{c02.eq.opt-bit} \\
    \text{\tpcifenglish{square bit}{Quadrat-Bit}}:\quad
      P_{\mathrm{opt}} &= \tfrac12\bigl(1 + \max(\abs{\delta},\, \abs{v_x},\, \abs{v_y})\bigr),
      \label{c02.eq.opt-quadrat} \\
    \text{Qubit}:\quad
      P_{\mathrm{opt}} &= \tfrac12\bigl(1 + \max(\abs{\delta},\, \norm{\vb{v}})\bigr).
      \label{c02.eq.opt-qubit}
  \end{align}
  Beim klassischen Bit ist das in den Koordinaten $x_1 = (\pw_0, \pw_1)$,
  $x_2 = (\pw'_0, \pw'_1)$ dasselbe wie
  \begin{equation}\label{c02.eq.opt-bit-p}
    P_{\mathrm{opt}} = \max(\lambda_1 \pw_0,\, \lambda_2 \pw'_0) + \max(\lambda_1 \pw_1,\, \lambda_2 \pw'_1) .
  \end{equation}
\end{proposition}
\begin{proof}
  Wir halten in Gleichung~\eqref{c02.eq.erfolg-koordinaten} zuerst $b$ fest und wählen
  $\vb{a}$ so gut wie möglich. Erlaubt sind alle $\vb{a}$, deren Länge höchstens
  $m = \min(b, 1-b)$ ist, wobei die Länge je nach Theorie $\abs{a}$,
  $\abs{a_x} + \abs{a_y}$ oder $\norm{\vb{a}}$ bedeutet. Der größte Wert von
  $\vb{v}\cdot\vb{a}$ ist dann $m\, L$, mit $L = \abs{v}$ beim Bit,
  $L = \max(\abs{v_x}, \abs{v_y})$ beim Quadrat-Bit und $L = \norm{\vb{v}}$ beim Qubit.
  Beim Bit wählt man dazu $a = \pm m$ mit dem Vorzeichen von $v$. Beim Quadrat-Bit legt man
  das ganze $m$ auf die Koordinate, in der $\vb{v}$ den größeren Betrag hat. Beim Qubit
  wählt man $\vb{a}$ in Richtung von $\vb{v}$.

  Es bleibt, $\lambda_2 + \delta\, b + L \min(b, 1-b)$ über $b \in [0,1]$ zu maximieren.
  Dieser Ausdruck ist auf $[0, \tfrac12]$ und auf $[\tfrac12, 1]$ jeweils eine Gerade in
  $b$; sein größter Wert liegt also bei $b = 0$, $b = \tfrac12$ oder $b = 1$. Die drei
  Werte sind $\lambda_2$, $\lambda_2 + \tfrac12\delta + \tfrac12 L = \tfrac12(1 + L)$ und
  $\lambda_2 + \delta = \lambda_1$. Wegen
  $\max(\lambda_1, \lambda_2) = \tfrac12(1 + \abs{\delta})$ ist der größte der drei
  $\tfrac12(1 + \max(\abs{\delta}, L))$.

  Zu \eqref{c02.eq.opt-bit-p}: Setzt man \eqref{c02.eq.bit-effekt} in
  \eqref{c02.eq.erfolg} ein, so erhält man $P = \lambda_2 + d_0 q_0 + d_1 q_1$ mit
  $d_0 = \lambda_1 \pw_0 - \lambda_2 \pw'_0$ und $d_1 = \lambda_1 \pw_1 - \lambda_2 \pw'_1$.
  Am größten wird das, wenn man $q_i = 1$ wählt, wo $d_i > 0$ ist, und sonst $q_i = 0$.
  Das ergibt $\lambda_2 \pw'_0 + \max(d_0, 0) + \lambda_2 \pw'_1 + \max(d_1, 0)$, und das ist
  die rechte Seite von \eqref{c02.eq.opt-bit-p}.
\end{proof}

Der Beweis liefert die optimalen Effekte gleich mit. Ist $\abs{\delta}$ die größte der
Zahlen in der Klammer, so ist Raten optimal: der Einseffekt, falls $\lambda_1$ die größere
Vorgabe ist, sonst die Nullfunktion. Andernfalls ist ein Extremaleffekt mit $b = \tfrac12$
optimal: beim Quadrat-Bit einer der vier Effekte $f_x^{\pm}$, $f_y^{\pm}$, beim Qubit der
Effekt $f_{\vb{n}}$ mit $\vb{n} = \vb{v}/\norm{\vb{v}}$. Die drei Formeln unterscheiden sich
nur darin, wie die Länge von $\vb{v}$ gemessen wird (Abbildung~\ref{c02.fig.laengen})\zoomref{c02.v4}.

\begin{tpcfigure}{c02.fig.laengen}
  \caption{Wie gut lässt sich ein Zustand $x_1$ vom Mittelpunkt $x_2$ des Zustandsraums
    unterscheiden, wenn beide gleich oft präpariert werden? Nach
    Proposition~\ref{c02.prop.optimum} ist dann $P_{\mathrm{opt}} = \tfrac12 + \tfrac14 L$,
    wobei $L$ die Länge des Ortsvektors von $x_1$ in der jeweiligen Theorie ist.
    Eingetragen sind die Orte von $x_1$ mit $P_{\mathrm{opt}} = \tfrac58$ und
    $P_{\mathrm{opt}} = \tfrac34$: (a) Punkte, (b) Quadrate, (c) Kreise. Mehr als
    $\tfrac34$ ist gegen den Mittelpunkt in keiner der drei Theorien zu erreichen.}
\end{tpcfigure}

\begin{vertiefung}[id=c02.v4,
    titel={Das Qubit-Optimum als Spurnorm},
    wozu={Gleichung~\eqref{c02.eq.opt-qubit} lässt sich ohne Koordinaten schreiben, direkt
      mit den Dichtematrizen der beiden Zustände.},
    voraussetzung={Abschnitt~\ref{c02.s4.3} und Proposition~\ref{c02.prop.optimum}.
      Der Kern braucht diesen Block nicht.},
    zurueck={Die Länge $\max(\abs{\delta}, \norm{\vb{v}})$ in
      Gleichung~\eqref{c02.eq.opt-qubit} ist eine Eigenschaft der Matrix
      $\lambda_1\rho_1 - \lambda_2\rho_2$ allein.}]
  Für eine hermitesche $2\times 2$-Matrix $\op{D}$ mit den Eigenwerten $\kappa_1$ und $\kappa_2$
  nennen wir $\norm{\op{D}}_1 = \abs{\kappa_1} + \abs{\kappa_2}$ ihre \emph{Spurnorm}.

  \begin{proposition}[Spurnormformel]\label{c02.prop.spurnorm}\claim[proof-included]{c02.cl.spurnorm}{derived}
    Seien $\rho_1 = \rho(\vb{r}_1)$ und $\rho_2 = \rho(\vb{r}_2)$ zwei Zustände des Qubits
    und $\lambda_1, \lambda_2 \geq 0$ mit $\lambda_1 + \lambda_2 = 1$. Dann ist
    \begin{equation}\label{c02.eq.spurnorm}
      P_{\mathrm{opt}} = \tfrac12\bigl(1 + \norm{\lambda_1 \rho_1 - \lambda_2 \rho_2}_1\bigr) .
    \end{equation}
  \end{proposition}
  \begin{proof}
    Nach Gleichung~\eqref{c02.eq.bloch} ist
    $\op{D} = \lambda_1\rho_1 - \lambda_2\rho_2
       = \tfrac12\bigl(\delta\, \Id + v_x \pauli{x} + v_y \pauli{y} + v_z \pauli{z}\bigr)$.
    Die Rechnung aus Proposition~\ref{c02.prop.bloch}(2) gibt für $\op{D}$ die Eigenwerte
    $\tfrac12(\delta + \norm{\vb{v}})$ und $\tfrac12(\delta - \norm{\vb{v}})$. Mit
    $L = \norm{\vb{v}}$ ist also
    $\norm{\op{D}}_1 = \tfrac12\bigl(\abs{\delta + L} + \abs{\delta - L}\bigr)$. Für
    $\delta \geq L$ ist das $\delta$, für $\delta \leq -L$ ist es $-\delta$, und für
    $-L \leq \delta \leq L$ ist es $L$; in jedem Fall also $\max(\abs{\delta}, L)$.
    Einsetzen in \eqref{c02.eq.opt-qubit} ergibt die Behauptung.
  \end{proof}

  Plávala beweist für jeden Zustandsraum eine Formel derselben Gestalt. An der Stelle der
  Spurnorm steht dort eine Norm auf $\AffDual{\StSp}$, die hier nicht eingeführt wird
  \srcref{plavala-2103.07469v2}{„Order unit norm, base norm and discrimination tasks“}.
\end{vertiefung}

Was folgt daraus für reine Zustände?

\begin{corollary}[Reine Zustände sind nicht immer sicher unterscheidbar]\label{c02.cor.rein}\claim[proof-included]{c02.cl.rein}{derived}
  \begin{enumerate}
    \item Im klassischen Bit und im Quadrat-Bit sind je zwei verschiedene reine Zustände
      sicher unterscheidbar.
    \item Zwei Zustände des Qubits sind genau dann sicher unterscheidbar, wenn beide rein
      sind und einander gegenüberliegen: $\vb{r}_2 = -\vb{r}_1$ mit $\norm{\vb{r}_1} = 1$.
    \item Die reinen Zustände $x_1 = (1,0,0)$ und $x_2 = (0,1,0)$ des Qubits sind
      verschieden und nicht sicher unterscheidbar. Für $\lambda_1 = \lambda_2 = \tfrac12$
      ist
      \begin{equation}\label{c02.eq.q1}
        P_{\mathrm{opt}} = \tfrac12\Bigl(1 + \tfrac12\sqrt{2}\Bigr) = \tfrac14\bigl(2 + \sqrt{2}\bigr) < 1 .
      \end{equation}
  \end{enumerate}
\end{corollary}
\begin{proof}
  (1) Beim Bit leistet der Effekt $(1,0)$ das Verlangte für $x_1 = s_0$, $x_2 = s_1$. Zwei
  verschiedene Ecken des Quadrats unterscheiden sich in mindestens einer Koordinate; diese
  ist in der einen Ecke $1$ und in der anderen $-1$. Der zugehörige Effekt $f_x^{\pm}$ oder
  $f_y^{\pm}$ hat dann die Werte $1$ und $0$.

  (2) Nach Proposition~\ref{c02.prop.perfekt} dürfen wir $\lambda_1 = \lambda_2 = \tfrac12$
  annehmen. Dann ist $\delta = 0$ und $\vb{v} = \tfrac12(\vb{r}_1 - \vb{r}_2)$, und
  Gleichung~\eqref{c02.eq.opt-qubit} gibt $P_{\mathrm{opt}} = 1$ genau dann, wenn
  $\norm{\vb{r}_1 - \vb{r}_2} = 2$ ist. Wegen
  $\norm{\vb{r}_1 - \vb{r}_2}^2 = \norm{\vb{r}_1}^2 + \norm{\vb{r}_2}^2 - 2\,\vb{r}_1\cdot\vb{r}_2 \leq 4$
  ist das nur möglich, wenn beide Längen $1$ sind und $\vb{r}_1\cdot\vb{r}_2 = -1$ ist,
  also $\vb{r}_2 = -\vb{r}_1$.

  (3) Hier ist $\vb{v} = \tfrac12(1, -1, 0)$ und $\norm{\vb{v}} = \tfrac12\sqrt{2}$.
\end{proof}

Der Zahlenwert in \eqref{c02.eq.q1} ist ungefähr \num[c02.cl.level-q1]{0{,}8536}{}. Teil (2)
gilt für die Kugel. Auf andere Zustandsräume lässt er sich nicht übertragen; im Quadrat sind
schon benachbarte Ecken sicher unterscheidbar\zoomref{c02.v5}.

Auch der umgekehrte Fall kommt vor: Die Zustände sind verschieden, und trotzdem ist Raten
das Beste. Für die Qubit-Zustände $x_1 = (0, 0, \tfrac14)$ und $x_2 = (0, 0, -\tfrac14)$
mit den Vorgaben $\lambda_1 = \tfrac{9}{10}$, $\lambda_2 = \tfrac{1}{10}$ ist
$\delta = \tfrac45$ und $\norm{\vb{v}} = \tfrac14$. Gleichung~\eqref{c02.eq.opt-qubit} gibt
$P_{\mathrm{opt}} = \num[c02.cl.level-q2]{\tfrac{9}{10}}{}$, und das erreicht schon der
Einseffekt: immer $x_1$ raten.

\begin{vertiefung}[id=c02.v5,
    titel={Vielecke und die Scheibe},
    wozu={Zwischen Quadrat und Kugel liegt eine ganze Familie von Zustandsräumen in der
      Ebene. An ihr sieht man, wie die Unterscheidbarkeit reiner Zustände von der Form
      des Zustandsraums abhängt, und was sich über das Verhältnis der beiden Kegel sagen
      lässt.},
    voraussetzung={Kern bis Korollar~\ref{c02.cor.rein}; für den letzten Teil
      Vertiefung~\ref{c02.v2}. Der Kern braucht diesen Block nicht.},
    zurueck={Ob sich zwei reine Zustände sicher unterscheiden lassen, hängt von der Form
      des Zustandsraums ab und nicht nur von den beiden Punkten.}]
  Für $n \geq 3$ sei $\StSp_n$ das regelmäßige $n$-Eck mit Umkreisradius $1$ und den Ecken
  \begin{equation}\label{c02.eq.vieleck}\claim{c02.cl.vieleck}{derived}
    \vb{v}_k = \bigl(\cos\tfrac{2\pi k}{n},\ \sin\tfrac{2\pi k}{n}\bigr), \quad k = 0, \dots, n-1,
    \qquad c_n = \cos\tfrac{\pi}{n} ,
  \end{equation}
  und $\StSp_{\mathrm{D}} = \{ \vb{r} \in \mathbb{R}^2 : \norm{\vb{r}} \leq 1 \}$ die
  Einheitsscheibe. Die reinen Zustände von $\StSp_n$ sind die Ecken $\vb{v}_k$, die von
  $\StSp_{\mathrm{D}}$ die Punkte des Randkreises. $\StSp_3$ ist ein Dreieck und damit ein
  Simplex, der Zustandsraum einer klassischen Theorie mit drei reinen Zuständen. $\StSp_4$
  ist ein Quadrat, dessen Ecken auf den Achsen liegen; aus $\StSp_{\mathrm{Q}}$ entsteht es
  durch Drehen und Verkleinern.

  \begin{proposition}[Extremaleffekte der Vielecke und der Scheibe]\label{c02.prop.vieleck}\claim[proof-omitted]{c02.cl.vieleck-effekte}{derived}
    Sei $\vb{m}_j = \bigl(\cos\tfrac{(2j+1)\pi}{n},\ \sin\tfrac{(2j+1)\pi}{n}\bigr)$ für
    $j = 0, \dots, n-1$, und seien
    \begin{equation}\label{c02.eq.vieleck-effekte}
      f_k(\vb{r}) = \frac{c_n + \vb{v}_k\cdot\vb{r}}{1 + c_n},
      \qquad
      g_j(\vb{r}) = \frac12 + \frac{\vb{m}_j\cdot\vb{r}}{2 c_n} .
    \end{equation}
    Die Extremaleffekte von $\StSp_n$ sind $0$ und $\UnitEff{\StSp}$ und außerdem: für
    ungerades $n$ die Effekte $f_k$ und $\UnitEff{\StSp} - f_k$ mit $k = 0, \dots, n-1$;
    für gerades $n$ die Effekte $g_j$ mit $j = 0, \dots, n-1$.
    Für ungerades $n$ sind das $2n + 2$ Extremaleffekte, für gerades $n$ sind es $n + 2$.
    Die Extremaleffekte von $\StSp_{\mathrm{D}}$ sind $0$, $\UnitEff{\StSp}$ und
    $\tfrac12(1 + \vb{n}\cdot\vb{r})$ für jeden Einheitsvektor $\vb{n}$ der Ebene.
  \end{proposition}

  Bei ungeradem $n$ hat $f_k$ den Wert $1$ in der Ecke $\vb{v}_k$ und den Wert $0$ auf der
  gegenüberliegenden Kante. Bei geradem $n$ hat $g_j$ den Wert $1$ auf einer Kante und den
  Wert $0$ auf der gegenüberliegenden. Für $n = 4$ sind die $g_j$ die vier Effekte
  $f_x^{\pm}$, $f_y^{\pm}$ des Quadrats in gedrehten Koordinaten. Die Herleitung steht bei
  den Nachweisen des Kapitels; für $n = 3, \dots, 8$ wurden die Listen zusätzlich durch
  Aufzählen aller Ecken der Effektmenge nachgerechnet.

  \emph{Unterscheiden.} Die Formel für die Scheibe ist dieselbe wie für die Kugel,
  Gleichung~\eqref{c02.eq.opt-qubit} mit Vektoren der Ebene. Für die Vielecke geht man die
  endlich vielen Extremaleffekte durch (Theorem~\ref{c02.thm.extremal}). Drei Ergebnisse,
  alle für $\lambda_1 = \lambda_2 = \tfrac12$:
  \begin{enumerate}
    \item Im Fünfeck sind die benachbarten reinen Zustände $\vb{v}_0$ und $\vb{v}_1$ nicht
      sicher unterscheidbar: Das Optimum ist
      $P_{\mathrm{opt}} = c_5 = \tfrac14(1 + \sqrt{5})$, ungefähr
      \num[c02.cl.level-p1]{0{,}8090}{}.
    \item Im Sechseck sind $\vb{v}_0 = (1,0)$ und
      $\vb{v}_2 = (-\tfrac12, \tfrac12\sqrt{3})$ reine Zustände und sicher unterscheidbar,
      $P_{\mathrm{opt}} = \num[c02.cl.level-h1]{1}{}$, obwohl sie einander nicht
      gegenüberliegen. Der Effekt $g_5$ hat in $\vb{v}_0$ den Wert $1$ und in $\vb{v}_2$
      den Wert $0$.
    \item Dasselbe Punktepaar $x_1 = (1,0)$, $x_2 = (-\tfrac34, \tfrac14)$ liegt im Quadrat
      $\StSp_4$, im Fünfeck $\StSp_5$ und in der Scheibe $\StSp_{\mathrm{D}}$
      (Abbildung~\ref{c02.fig.vielecke}). Die Optima sind
      \begin{equation}\label{c02.eq.gleiches-paar}\claim{c02.cl.gleiches-paar}{derived}
        P_{\mathrm{opt}}(\StSp_4) = 1, \qquad
        P_{\mathrm{opt}}(\StSp_5) = \frac{55 - 7\sqrt{5}}{40}, \qquad
        P_{\mathrm{opt}}(\StSp_{\mathrm{D}}) = \frac{8 + 5\sqrt{2}}{16} ,
      \end{equation}
      ungefähr \num[c02.cl.level-c1]{1}{}, \num[c02.cl.level-c2]{0{,}9837}{} und
      \num[c02.cl.level-c3]{0{,}9419}{}.
  \end{enumerate}
  Die dritte Aussage hängt daran, wie die drei Zustandsräume zueinander liegen. Mit dem
  achsenparallelen Quadrat $\StSp_{\mathrm{Q}} = [-1,1]^2$ und der Einheitsscheibe in
  denselben Koordinaten gibt es kein solches Paar: Jedes Paar von Punkten der Scheibe, das
  in $\StSp_{\mathrm{Q}}$ sicher unterscheidbar ist, besteht aus zwei gegenüberliegenden
  Randpunkten der Scheibe und ist auch dort sicher unterscheidbar.

  \begin{tpcfigure}{c02.fig.vielecke}
    \caption{Dasselbe Punktepaar $x_1 = (1,0)$, $x_2 = (-\tfrac34, \tfrac14)$ in drei
      Zustandsräumen, jeweils mit den Niveaulinien eines optimalen Extremaleffekts. (a) Im
      Quadrat $\StSp_4$ liegt $x_2$ auf der Kante, auf der der Effekt den Wert $0$ hat; die
      Unterscheidung gelingt sicher. (b) Im Fünfeck liegt $x_2$ im Inneren, knapp neben
      der Kante mit $f = 0$. (c) In der Scheibe berührt die Linie $f = 0$ den Rand in einem
      einzigen Punkt, und $x_2$ liegt deutlich daneben. Die Optima stehen in
      \eqref{c02.eq.gleiches-paar}.}
  \end{tpcfigure}

  \emph{Die beiden Kegel.} In den Koordinaten von Vertiefung~\ref{c02.v2} gehört zum
  $n$-Eck mit Umkreisradius $\gamma > 0$ der Zustandskegel
  $C_{n,\gamma} = \{ t\,(1, \gamma\,\vb{r}) : t \geq 0,\ \vb{r} \in \StSp_n \}$.

  \begin{proposition}[Wann die Kegel der Vielecke gleich sind]\label{c02.prop.selbstdual}\claim[proof-omitted]{c02.cl.selbstdual}{derived}
    Sei $C_{n,\gamma}^{*} = \{ (s, \vb{a}) : s\,t + \vb{a}\cdot\vb{R} \geq 0
    \ \text{\tpcifenglish{for all}{für alle}}\ (t, \vb{R}) \in C_{n,\gamma} \}$ der duale
    Kegel bezüglich des gewöhnlichen Skalarprodukts des $\mathbb{R}^3$.
    \begin{enumerate}
      \item Für $\gamma = 1$ ist $C_{n,1}^{*} \neq C_{n,1}$ für jedes $n \geq 3$.
      \item Für ungerades $n$ ist $C_{n,\gamma}^{*} = C_{n,\gamma}$ genau dann, wenn
        $\gamma^2 = 1/c_n$ ist.
      \item Für gerades $n$ ist $C_{n,\gamma}^{*} \neq C_{n,\gamma}$ für jedes
        $\gamma > 0$.
      \item Für die Scheibe ist der Kegel $\{ (t, \vb{R}) : t \geq \norm{\vb{R}} \}$ gleich
        seinem dualen Kegel.
    \end{enumerate}
  \end{proposition}

  Die Aussagen (1) bis (3) wurden für $n = 3, \dots, 8$ zusätzlich numerisch geprüft. Sie
  betreffen die Gleichheit zweier Mengen in festen Koordinaten mit dem gewöhnlichen
  Skalarprodukt. Für jedes $n$ und jedes $\gamma$ gibt es eine umkehrbare lineare
  Abbildung, die $C_{n,\gamma}^{*}$ auf $C_{n,\gamma}$ abbildet, eine Drehung verbunden
  mit einer Streckung; das ist eine schwächere Aussage als Gleichheit. Ob sich für gerades
  $n$ ein anderes Skalarprodukt finden lässt, für das die Kegel gleich werden, ist hier
  nicht untersucht. Auch das Dreieck fällt unter (1): In den Wahrscheinlichkeitskoordinaten
  $(\pw_1, \pw_2, \pw_3)$ ist der Zustandskegel des Simplex gleich seinem dualen Kegel, in
  den Koordinaten des regelmäßigen Dreiecks mit Umkreisradius $1$ nicht. Die beiden
  Koordinatensysteme gehen nicht durch eine Drehung auseinander hervor.
\end{vertiefung}

\section{Das Spiel „Distinguish“}\label{c02.s7}

Zu diesem Kapitel gehört ein Spiel. Diese Seite ist seine Anleitung.

\emph{Regeln.} Ein Level gibt eine Theorie vor, zwei Zustände $x_1$, $x_2$ und zwei
Vorgaben $\lambda_1$, $\lambda_2$. Gespielt wird die Aufgabe aus Abschnitt~\ref{c02.s6}: Man
zieht einen Effekt $f$ an eine Stelle der Effektmenge. Gemessen wird mit dem Paar $f$,
$\UnitEff{\StSp} - f$; bei Ja wird $x_1$ geraten, bei Nein $x_2$. Es gibt keinen Zufall
und keine Zeitgrenze.

\emph{Die Erfolgsgröße} ist $P(f)$ aus Gleichung~\eqref{c02.eq.erfolg}, die
Wahrscheinlichkeit, richtig zu raten. Sie ist die einzige physikalische Größe des Spiels. Das Spiel
zeigt sie für den gewählten Effekt an.

\emph{Die Belohnung} wird aus $P(f)$ und $P_{\mathrm{opt}}$ berechnet und ist eine
Entscheidung der Spielgestaltung, keine Physik. Drei Sterne gibt es nur für
$P(f) = P_{\mathrm{opt}}$; verglichen wird exakt, ohne Toleranz. Darunter gibt es zwei
Stufen mit festen Schwellen $0 < \varepsilon_2 < \varepsilon_1$: zwei Sterne für
$P_{\mathrm{opt}} - P(f) \leq \varepsilon_2$, sonst einen Stern für
$P_{\mathrm{opt}} - P(f) \leq \varepsilon_1$, sonst keinen. Die beiden Schwellen sind
gewählte Entwurfswerte, nicht hergeleitet und nicht geprüft; ihre Werte stehen in der
Regeltafel des Spiels und nicht in diesem Buch.

\emph{Das bewiesene Optimum} $P_{\mathrm{opt}}$ jedes Levels steht in
Tabelle~\ref{c02.tab.level}. Es folgt aus Proposition~\ref{c02.prop.optimum} und wird nach
Theorem~\ref{c02.thm.extremal} von einem Extremaleffekt erreicht. Mehr ist in keinem Level
möglich.

\begin{tpctable}{c02.tab.level}
  \caption{Die Level der ersten Ausgabe. Bit: Zustände als $(\pw_0, \pw_1)$, Effekte als
    $(q_0, q_1)$. Quadrat-Bit und Qubit: Zustände als $\vb{r}$. Letzte Spalte: ein
    optimaler Extremaleffekt; in Q1 ist $\vb{n} = (1, -1, 0)/\sqrt{2}$.}
  \begin{tabular}{llllll}
    \toprule
    Level & $x_1$ & $x_2$ & $(\lambda_1, \lambda_2)$ & $P_{\mathrm{opt}}$ & optimal \\
    \midrule
    B1 & $(1, 0)$ & $(0, 1)$ & $(\tfrac12, \tfrac12)$ & \num[c02.cl.level-b1]{1}{} & $(1, 0)$ \\
    B2 & $(\tfrac34, \tfrac14)$ & $(\tfrac14, \tfrac34)$ & $(\tfrac12, \tfrac12)$ & \num[c02.cl.level-b2]{\tfrac34}{} & $(1, 0)$ \\
    S1 & $(1, 0)$ & $(-1, 0)$ & $(\tfrac12, \tfrac12)$ & \num[c02.cl.level-s1]{1}{} & $f_x^{+}$ \\
    S2 & $(1, 1)$ & $(-1, -1)$ & $(\tfrac12, \tfrac12)$ & \num[c02.cl.level-s2]{1}{} & $f_x^{+}$, $f_y^{+}$ \\
    Q1 & $(1, 0, 0)$ & $(0, 1, 0)$ & $(\tfrac12, \tfrac12)$ & \num[c02.cl.level-q1]{\tfrac14(2 + \sqrt{2})}{} & $f_{\vb{n}}$ \\
    Q2 & $(0, 0, \tfrac14)$ & $(0, 0, -\tfrac14)$ & $(\tfrac{9}{10}, \tfrac{1}{10})$ & \num[c02.cl.level-q2]{\tfrac{9}{10}}{} & $\UnitEff{\StSp}$ \\
    \bottomrule
  \end{tabular}
\end{tpctable}

\begin{tpcfigure}{c02.fig.unterscheiden}
  \demo{c02.demo.distinguish}
  \caption{Je ein Level in den drei Theorien mit den Niveaulinien eines optimalen
    Extremaleffekts (gestrichelt). (a) B2: Der Effekt $(1,0)$ hat in $x_1$ den Wert
    $\tfrac34$ und in $x_2$ den Wert $\tfrac14$. (b) S2: Der Effekt $f_x^{+}$ trennt die
    beiden Ecken sicher; $f_y^{+}$ auch. (c) Q1, Schnitt $r_z = 0$ durch die Zustandskugel des Qubits:
    Keine Schar paralleler Linien hat $x_1$ auf $f = 1$ und zugleich $x_2$ auf $f = 0$.}
\end{tpcfigure}

In S1 sind beide Zustände gemischt und trotzdem sicher unterscheidbar. In S2 sind zwei
Extremaleffekte optimal und mit ihnen die ganze Kante dazwischen
(Abbildung~\ref{c02.fig.optimum}(b)). In Q2 ist Raten optimal: Der beste Effekt ist der
Einseffekt, immer $x_1$ raten.

\section{Von hier zur Quantentheorie}\label{c02.s8}

Der Rahmen dieses Kapitels zeichnet keinen der drei Zustandsräume aus: Strecke, Quadrat und
Kugel erfüllen alle die Bedingungen aus \defref{D3}. Welcher Zustandsraum zu einem
gegebenen System gehört, sagt der Rahmen nicht. In der Quantentheorie treten die
Zustandsräume einer bestimmten Familie auf, zu der die Kugel aus
Abschnitt~\ref{c02.s4.3} als kleinstes nichttriviales Beispiel gehört. Zum Schritt dorthin
steht in der Rahmung des Kurses der Satz:
% GAP-TJO: Wie TJO den Schritt von den GPTs zur Quantentheorie benannt haben will, ist
% offen (evidence/gaps.json, G-0012).  Die Bezeichnung „Modellwahl“ aus
% docs/plan/PLAN.md, Abschnitt 13.1, A11 ist nicht sein Wortlaut und steht deshalb nicht
% im Text; sein Satz ist der folgende (BRIEF.md:47).  Die deutsche Fassung ist eine
% Übersetzung.
\begin{equation}\label{c02.eq.rahmen-spezialisierung}\claim{c02.cl.rahmen-spezialisierung}{author-stipulated}
  \begin{aligned}
    &\text{\tpcifenglish{The GPT presentation allows us to present}{Die GPT-Darstellung erlaubt uns, die Postulate}} \\
    &\text{\tpcifenglish{the postulates of a physical theory which we}{einer physikalischen Theorie anzugeben, die wir}} \\
    &\text{\tpcifenglish{subsequently specialise to QM next.}{anschließend zur Quantenmechanik spezialisieren.}}
  \end{aligned}
\end{equation}

\begin{fundament}[id=c02.f1,
    titel={Ordnungseinheitsräume},
    wozu={Die Effekte eines Zustandsraums bilden ein Intervall in einem geordneten
      Vektorraum mit Ordnungseinheit. Hier steht diese Struktur so, wie Plávala sie für
      endliche Dimension angibt, und es wird gesagt, was der Kurs darüber hinaus noch
      schuldet.},
    voraussetzung={Vertiefung~\ref{c02.v2}; Halbordnungen. Weder Kern noch Vertiefungen
      brauchen diesen Block.},
    zurueck={Man kann eine Theorie mit ihren Zuständen beginnen oder mit ihren Effekten;
      in endlicher Dimension führt beides zu demselben Rahmen.}]
  Auf $\Aff{\StSp}$ gibt es eine natürliche Ordnung: $f \leq g$ bedeutet
  $f(x) \leq g(x)$ für alle $x \in \StSp$, also $g - f \in \AffPos{\StSp}$. Ein
  \emph{geordneter Vektorraum} ist bei Plávala ein Vektorraum mit einer Relation $\leq$,
  die reflexiv, antisymmetrisch und transitiv ist und sich mit der Addition und mit der
  Multiplikation mit nicht negativen Zahlen verträgt
  \srcref{plavala-2103.07469v2}{Anhang „Convex cones and ordered vector spaces“}.

  \tpcdef{D6}

  \begin{proposition}[Affine Funktionen als Ordnungseinheitsraum]\label{c02.prop.oeraum}\claim[proof-omitted]{c02.cl.oeraum}{source-backed}
    Für jeden Zustandsraum $\StSp$ ist $(\Aff{\StSp}, \leq, \UnitEff{\StSp})$ ein
    Ordnungseinheitsraum, und die Effekte sind darin das Ordnungsintervall zwischen $0$
    und der Ordnungseinheit:
    \begin{equation}\label{c02.eq.intervall}
      \EffAlg{\StSp} = \{ f \in \Aff{\StSp} : 0 \leq f \leq \UnitEff{\StSp} \} .
    \end{equation}
  \end{proposition}

  Plávala beweist die drei Bestandteile einzeln: dass $\Aff{\StSp}$ mit dieser Ordnung ein
  geordneter Vektorraum ist, dass $\UnitEff{\StSp}$ eine Ordnungseinheit ist, und
  Gleichung~\eqref{c02.eq.intervall}. Quelle:
  \srcref{plavala-2103.07469v2}{„Effect algebra“}.
  Abbildung~\ref{c02.fig.ordnungsintervall} zeigt das Intervall für das klassische Bit.

  \begin{tpcfigure}{c02.fig.ordnungsintervall}
    \caption{Die Effektmenge des klassischen Bits als Ordnungsintervall, in den Koordinaten
      $(q_0, q_1)$ aus Proposition~\ref{c02.prop.bit}. Der positive Kegel $\AffPos{\StSp}$
      ist der Quadrant rechts oben vom Nullpunkt. Die Bedingung $f \leq \UnitEff{\StSp}$
      sagt, dass $\UnitEff{\StSp} - f$ in diesem Kegel liegt; das ist der Quadrant links
      unten von $\UnitEff{\StSp}$. Der Durchschnitt der beiden Quadranten ist das
      Einheitsquadrat $\EffAlg{\StSp}$.}
  \end{tpcfigure}

  Der Weg geht auch in die andere Richtung, von den Effekten zu den Zuständen. Ein Kegel
  $C$ in einem Vektorraum $V$ heißt dabei \emph{spitz}, wenn außer $0$ kein Vektor zugleich
  mit seinem Negativen in $C$ liegt, und \emph{erzeugend}, wenn $C$ den Raum $V$
  aufspannt.

  \begin{proposition}[Von den Effekten zurück zu den Zuständen]\label{c02.prop.rueckweg}\claim[proof-omitted]{c02.cl.rueckweg}{source-backed}
    Sei $V$ ein reeller, endlichdimensionaler Vektorraum, $C \subset V$ ein konvexer,
    abgeschlossener, spitzer und erzeugender Kegel und $\OrdUnit \in C$ eine
    Ordnungseinheit. Sei $E = \{ v \in V : 0 \leq v \leq \OrdUnit \}$ und
    \begin{equation}\label{c02.eq.rueckweg}
      S(E) = \{ \psi \in C^{*} : \pair{\psi}{\OrdUnit} = 1 \} ,
    \end{equation}
    wobei $C^{*}$ der duale Kegel von $C$ ist. Dann ist $S(E)$ ein Zustandsraum, und seine
    Effektmenge ist wieder $E$: Es gilt $\EffAlg{S(E)} = E$.
  \end{proposition}

  Plávala gibt dafür die Beweisidee an und folgert, dass der Rahmen, der von solchen
  Intervallen ausgeht, derselbe ist wie der, der von Zustandsräumen ausgeht
  \srcref{plavala-2103.07469v2}{„Connection to abstract convex effect algebras and order unit spaces“}.
  Dort steht auch der abstrakte Begriff der Effektalgebra, den dieses Kapitel nicht
  braucht.

  \emph{Was offen bleibt.} In der Rahmung des Kurses steht an dieser Stelle der Satz
  % GAP-TJO: BRIEF.md:47 verlangt „Archimedean order unit spaces“.  Das Wort „Archimedean“
  % kommt in der Quelle (Plávala v2) nicht vor, und keine Quelle auf der Platte definiert
  % es (evidence/gaps.json, G-0002).  Was es für den Kurs bedeutet, wie es sich zu
  % Definition D6 verhält und wie der Satz in unendlicher Dimension zu lesen ist, muss TJO
  % sagen.  Die deutsche Fassung des Satzes ist eine Übersetzung.
  \begin{equation}\label{c02.eq.rahmen-archimedisch}\claim{c02.cl.rahmen-archimedisch}{author-stipulated}
    \begin{aligned}
      &\text{\tpcifenglish{All physical theories, for me,}{Alle physikalischen Theorien sind für mich}} \\
      &\text{\tpcifenglish{are Archimedean order unit spaces.}{archimedische Ordnungseinheitsräume.}}
    \end{aligned}
  \end{equation}
  Dieses Kapitel deckt davon nur einen Teil ab. Erstens ist \defref{D6} die Definition aus
  der Quelle, und in ihr kommt die Eigenschaft „archimedisch“ nicht vor; was sie bedeutet
  und warum der Kurs sie verlangt, ist hier nicht gesagt. Zweitens sind alle Aussagen
  dieses Kapitels unter der Bedingung (S5) bewiesen oder zitiert, also für
  endlichdimensionale Zustandsräume. Über unendlichdimensionale Theorien sagt es nichts.
\end{fundament}
